\documentclass[11pt,a4paper,titlepage]{ltjsbook}

\usepackage{luatexja}
\usepackage{amsmath,amssymb,amsfonts,amsthm}
\usepackage{geometry}
\geometry{top=23mm,bottom=23mm,left=23mm,right=23mm}
\usepackage{tikz}
\usetikzlibrary{arrows.meta,calc,positioning,decorations.pathmorphing,shapes.geometric,backgrounds}
\usepackage{tcolorbox}
\tcbuselibrary{skins,breakable}
\usepackage{hyperref}
\usepackage{booktabs}
\usepackage{enumitem}

\hypersetup{
    colorlinks=true,
    linkcolor=blue!70!black,
    citecolor=green!50!black,
    urlcolor=teal!80!black,
    pdftitle={古典で学び直すアインシュタイン相対性理論入門},
    pdfauthor={Albert Einstein / 著・解説}
}

% Custom boxes
\newtcolorbox{lecturebox}[1][大人のための徹底解説講義]{
    enhanced,
    colback=blue!5!white,
    colframe=blue!65!black,
    fonttitle=\bfseries\gtfamily,
    title={#1},
    arc=2mm,
    breakable,
    boxrule=1pt,
    top=3mm, bottom=3mm, left=3mm, right=3mm
}

\newtcolorbox{origbox}[1][アインシュタインの原典思索]{
    enhanced,
    colback=yellow!5!white,
    colframe=orange!70!black,
    fonttitle=\bfseries\gtfamily,
    title={#1},
    arc=2mm,
    breakable,
    boxrule=1pt,
    top=3mm, bottom=3mm, left=3mm, right=3mm
}

\newtcolorbox{examplebox}[1]{
    enhanced,
    colback=green!5!white,
    colframe=teal!70!black,
    fonttitle=\bfseries\gtfamily,
    title={#1},
    arc=2mm,
    breakable,
    boxrule=1pt,
    top=3mm, bottom=3mm, left=3mm, right=3mm
}

\newtcolorbox{exercisebox}[1]{
    enhanced,
    colback=purple!5!white,
    colframe=purple!70!black,
    fonttitle=\bfseries\gtfamily,
    title={#1},
    arc=2mm,
    breakable,
    boxrule=1pt,
    top=3mm, bottom=3mm, left=3mm, right=3mm
}

\title{
    {\Large\gtfamily 数学・科学古典名著再学習シリーズ 第6弾}\\[8pt]
    {\Huge\bfseries\gtfamily 古典で学び直す\\アインシュタイン相対性理論入門}\\[10pt]
    {\large\gtfamily いまからでも遅くない、高校までの数学で理解する歴史的名著}\\[6pt]
    {\normalsize 原典：Albert Einstein \textit{Relativity: The Special and the General Theory} (1916/1920)}
}
\author{アルベルト・アインシュタイン 著\\編訳・現代的解説：プロジェクトGutenberg名著復刊委員会}
\date{2026年10月 完全決定版}

\begin{document}

\maketitle

\frontmatter

\chapter*{はじめに：大人の学び直しとしての『相対性理論』}
\addcontentsline{toc}{chapter}{はじめに：大人の学び直しとしての『相対性理論』}

1916年、アルベルト・アインシュタイン（Albert Einstein, 1879--1955）は、自ら打ち立てた二大物理学革命――特殊相対性理論（1905年）と一般相対性理論（1915年）――の全容を、物理学者だけでなく\textbf{「科学的・哲学的な興味を持つ一般読者」}に向けて自ら平易に解説した稀代の名著『特殊および一般相対性理論について（原題：Über die spezielle und die allgemeine Relativitätstheorie, gemeinverständlich）』を著しました。

世間一般において「相対性理論」といえば、難解極まるテンソル解析や微分幾何学、あるいはSFの世界のタイムトラベルやブラックホールといった先入観が先行しがちです。しかし、アインシュタイン自身が序文で宣言しているように、本書の目的は\textbf{「大学の専門教育を受けていない読者でも、標準的な高校程度の教育（中等教育）を受けていれば、その核心思想と論理を完全に追体験できるようにすること」}にありました。

本入門・解説版（第6弾）では、人生経験豊かな50代〜80代の大人の学び直し読者のために、以下の三大特色を備えて構成されています：
\begin{enumerate}[label=\textbf{\arabic*.}]
    \item \textbf{アインシュタイン自身の思考実験を高校数学で完全追体験}：列車と堤防の思考実験、光時計、動く物体の縮み、同時性の不一致、重力と加速系の等価原理を、三平方の定理と1次方程式だけで完全に導出します。
    \item \textbf{全12章・高精細TikZ図版12点と実社会の応用例}：GPS衛星のナノ秒単位の時計補正、新幹線・ジェット機実験、素粒子ミューオンの寿命延伸、巨大重力レンズ、宇宙膨張と現代ビッグバン宇宙論など、現代社会と科学の礎を実感できます。
    \item \textbf{重要用語辞典＆総合演習問題と詳細解答}：基礎から段階的に確認できる総合問題（全10問）とステップ・バイ・ステップの模範解答を収録し、独学でも完全に理解を深められる構成となっています。
\end{enumerate}

知的好奇心に満ちた物理学の最高峰へ、アインシュタインと共に歩みを進めましょう。

\clearpage
\tableofcontents
\clearpage

\mainmatter


\clearpage
\part{特殊相対性理論（時間と空間の革命）}

\chapter{幾何学の物理的意味と座標系}
\begin{center}
{\Large\bfseries\gtfamily ユークリッド幾何学の真実性と剛体定規による空間測定の再定義}
\end{center}
\vspace*{2mm}

\section*{本章の視覚的見取り図}
\begin{center}
\begin{tikzpicture}[scale=1.0, >=Stealth]
    % Coordinates and grid
    \draw[->, thick, gray!60] (-0.5,0) -- (6,0) node[right, black] {$x$ (空間)};
    \draw[->, thick, gray!60] (0,-0.5) -- (0,4) node[above, black] {$y$};
    \draw[step=1.0, gray!20, very thin] (-0.5,-0.5) grid (5.5,3.5);
    % Rigid body / Reference frame
    \node[draw, rectangle, rounded corners, fill=blue!10, draw=blue!70, thick, inner sep=4pt] at (1.5,2.5) {基準座標系 $K$};
    % Measuring rod
    \draw[|-|, ultra thick, red!80!black] (1,0.5) -- (4,0.5) node[midway, above, font=\small\bfseries] {基準剛体定規 ($L_0$)};
    % Event points
    \fill[red] (1,0.5) circle (2.5pt) node[below=2pt, font=\small] {$A (x_1, y_1)$};
    \fill[red] (4,0.5) circle (2.5pt) node[below=2pt, font=\small] {$B (x_2, y_2)$};
    % Annotation
    \node[align=left, fill=white, draw=gray!40, rounded corners, font=\footnotesize, inner sep=3pt] at (3.5, 2.2) {幾何学の命題：\\「2点間の距離は定規の\\目盛り数として実測される」};
\end{tikzpicture}

\small\textbf{図版解説：}ユークリッド幾何学の命題を物理的な「剛体定規」による目盛り計測として捉え直す座標系 $K$ の構造。
\end{center}
\vspace*{2mm}

\section*{アインシュタインの原典思索}
\begin{origbox}[原著第1章：アルベルト・アインシュタインの肉声]
幾何学は特定の公理（たとえば「2点を通る直線はただ1つ存在する」など）から論理的に命題を演繹する体系である。しかし純粋幾何学においては、「直線」や「点」という概念が現実の経験的対象とどう対応するかという問いには答えない。
この幾何学的命題に物理的内容を与えるためには、たとえば「直線」とは「光線の軌跡」あるいは「真っ直ぐな剛体定規の縁」であるというように、現実の物理的実体との対応関係（規約）を定めなければならない。このとき幾何学は自然科学の最も古い一分野となる。
\end{origbox}

\section*{大人のための徹底解説講義}
\begin{lecturebox}[徹底講義：高校数学で追体験する物理学革命]
私たちが中学や高校で学んだユークリッド幾何学は、公理から出発して定理を純粋論理的に証明していく美しい学問でした。しかし、アインシュタインは冒頭で極めて根本的な問いを投げかけます。
「私たちがノートの上に引いた直線や、教科書に書かれた三角形の内角の和が180度であるという命題は、現実の物理的宇宙においても絶対に真実なのだろうか？」

純粋数学としての幾何学は、「点」や「直線」が現実の何であるかを定義しません。しかし、大工さんが家を建てるとき、測量士が土地の境界を定めるとき、私たちは金属や木で作られた「定規（剛体）」を現実に並べて長さを測ります。
アインシュタインは、「幾何学の命題とは、剛体定規を実際に当てがって測定した結果についての物理的主張である」と明確に位置づけ直しました。そして、物体の位置を一意に指定するために「直交座標系（デカルト座標系 $x, y, z$）」を導入します。
空間とは単なる抽象的な入れ物ではなく、現実の測定器（定規）によって目盛りが振られた「物理的座標系」なのです。
\end{lecturebox}

\vspace*{3mm}
\begin{examplebox}{大人の学び直しのための実社会・重要例題：第1章}
\textbf{【実社会の具体例：土地の境界確定とGPS・電子基準点の幾何学】}\\
私たちが家や土地を売買する際、法務局に登記される「公図」や「地積測量図」は、一見すると何千年ものあいだ変わらないユークリッド幾何学の平面上に描かれています。\\
(1) アインシュタインが指摘した「幾何学の命題（直線や距離）は、物理的な剛体定規を当てがうことによって初めて物理的意味を持つ」という主張を、現代の国土地理院による電子基準点網（GNSS測量）と対比して説明せよ。\\
(2) 地球が球体（曲面）であるにもかかわらず、なぜ日常生活の建築や不動産登記では「ユークリッド幾何学」を何ら問題なく信用して適用できるのか、微小領域における平坦性の観点から考察せよ。
\\[2mm]
\textbf{【懇切丁寧な解答・ステップ解説】}\\
\textbf{(1) の解：}
純粋数学におけるユークリッド幾何学の「点」や「直線」は、頭の中の完全無欠な抽象概念であり、それ自体では現実世界のどこにも存在しません。アインシュタインは、この抽象的な命題に「金属の定規（剛体）を並べて長さを測る」という物理的実操作を割り当てて初めて、幾何学が自然科学（物理学）の一部になると喝破しました。
現代の電子基準点測量でも、衛星から届く「電波の到達時間（物理量）」に光速を乗じて距離を算出し、三角網を構築しています。幾何学は机上の空論ではなく、常に測定手段という物理的実体と分かちがたく結びついているのです。
\\[2mm]
\textbf{(2) の解：}
地球全体の半径は約 $6371\,\text{km}$ という長大さを持っています。私たちが生活する数百メートル四方の住宅地や建築現場では、球面による曲がり（ズレ）はミリメートル以下という無視できる微小量にとどまります。数学的に言えば、「いかに複雑に湾曲した多様体であっても、局所的（ごく狭い範囲）には平坦なユークリッド空間と見なせる」という微積分・幾何学の基本原理によって、私たちは安心して平面幾何学を実用できているのです。
\end{examplebox}


\clearpage
\chapter{古典力学における空間と時間}
\begin{center}
{\Large\bfseries\gtfamily ガリレイ変換の限界と「光速度不変の原理」の登場}
\end{center}
\vspace*{2mm}

\section*{本章の視覚的見取り図}
\begin{center}
\begin{tikzpicture}[scale=0.9, >=Stealth]
    % Embankment
    \draw[thick, fill=brown!20] (-1,-0.5) rectangle (7,0.2);
    \node at (3,-0.15) {\textbf{堤防（静止座標系 $K$）}};
    \node[draw, circle, fill=yellow!40, inner sep=2pt] (obsK) at (0.5, 0.7) {観測者 $M$};
    % Train
    \draw[thick, fill=blue!15, rounded corners] (1.5,1.2) rectangle (6.5,2.6);
    \node at (4,2.2) {\textbf{走行列車（運動系 $K'$）}};
    \draw[->, ultra thick, red!80!black] (6.5, 1.9) -- (7.8, 1.9) node[right] {速度 $v$};
    \node[draw, circle, fill=cyan!40, inner sep=2pt] (obsKp) at (3.5, 1.7) {観測者 $M'$};
    % Walker inside train
    \draw[->, thick, teal] (4.5, 1.6) -- (5.5, 1.6) node[midway, above, font=\scriptsize] {歩行 $w$};
    % Classical addition
    \node[align=left, fill=white, draw=gray!50, rounded corners, font=\footnotesize] at (3, 3.4) {ガリレイ変換の古典則：$W = v + w$\\しかし「光速 $c$」に対してこの単純加算は通用するのか？};
\end{tikzpicture}

\small\textbf{図版解説：}堤防の静止系 $K$ と一定速度 $v$ で走る列車系 $K'$ における運動の相対性と古典速度合成則。
\end{center}
\vspace*{2mm}

\section*{アインシュタインの原典思索}
\begin{origbox}[原著第2章：アルベルト・アインシュタインの肉声]
古典力学はガリレイ・ニュートンの相対性原理に立脚している。すなわち、ある基準座標系 $K$ において力学法則が成立するならば、これに対して等速直線運動を行っている別の座標系 $K'$ においても全く同一の力学法則が成立する。
線路の堤防に対して速度 $v$ で走る列車の中を、乗客が列車の進行方向に速度 $w$ で歩くとき、堤防に対する乗客の速度は古典的には $W = v + w$ となる。しかし、この常識的な加算則を光の伝播（速度 $c$）に適用しようとすると、マクスウェル電磁気学の基本法則と激しく衝突することになる。
\end{origbox}

\section*{大人のための徹底解説講義}
\begin{lecturebox}[徹底講義：高校数学で追体験する物理学革命]
時速 $100\,\text{km}$ で走る電車の中で、前方へ時速 $10\,\text{km}$ でボールを投げると、地上に立っている人から見ればボールは時速 $110\,\text{km}$ で飛んでいきます。これが高校物理で習う「ガリレイの速度加算則（$W = v + w$）」であり、ニュートン以来200年以上にわたって人類の絶対的な常識でした。

ところが19世紀末、マクスウェルの電磁気学が完成すると、驚くべき矛盾が浮上しました。真空中を進む電磁波（光）の速度 $c \approx 30万\,\text{km/s}$ は、光源が動いていようが止まっていようが、常に一定の値をとると方程式が示していたのです。
もしガリレイの法則が正しければ、光速の $90\%$ で飛ぶロケットから前方に放たれた光は、地球から見れば「光速の $1.9$ 倍」に見えなければなりません。しかし、マイケルソン・モーリーの実験をはじめとするあらゆる精密測定は、「光の速度は誰から見ても秒速30万キロメートルのままであり、絶対に変わらない」という奇妙な事実を突きつけました。
アインシュタインは、この一見両立し得ない2つの原理――「相対性原理」と「光速度不変の原理」――を両方とも真理として受け入れる決断を下します。
\end{lecturebox}

\vspace*{3mm}
\begin{examplebox}{大人の学び直しのための実社会・重要例題：第2章}
\textbf{【実社会の具体例：空港の動く歩道と「時速 $100\,\text{km}$ の列車から光を投げたら？」】}\\
空港の長い通路にある「動く歩道（時速 $4\,\text{km}$）」の上を、旅行者が前方へ時速 $4\,\text{km}$ で歩くと、外の床に立っている人から見れば速度は $4 + 4 = 8\,\text{km/h}$ に見えます。これは高校物理で習う「ガリレイの速度合成則」です。\\
(1) もし時速 $100\,\text{km}$ で走る列車のヘッドライトから放たれた光の速度を、線路端の静止観察者が測ったとき、ニュートン力学の常識では「光速 $c + 100\,\text{km/h}$」になるはずである。なぜこの結論がマクスウェルの電磁気学と衝突したのかを述べよ。\\
(2) アインシュタインが、ニュートン力学のガリレイ変換を捨てて「光速度不変の原理」を公理として採択した決断の革命性を考察せよ。
\\[2mm]
\textbf{【懇切丁寧な解答・ステップ解説】}\\
\textbf{(1) の解：}
19世紀末に確立したマクスウェルの電磁気学によれば、光は電磁波の一種であり、その真空中での伝播速度 $c$ は真空の誘電率 $\varepsilon_0$ と透磁率 $\mu_0$ によって $c = 1/\sqrt{\varepsilon_0 \mu_0} \approx 3 \times 10^8\,\text{m/s}$ と一意に定まります。ここには「光源の速度」というパラメータは一切含まれていません。
もしガリレイ変換が正しいなら、電磁気学の方程式は動いている列車の中では成り立たなくなってしまいます。「力学の常識（ガリレイ変換）」と「電磁気学の基本法則」は真っ向から衝突していたのです。
\\[2mm]
\textbf{(2) の解：}
当時の多くの物理学者は「電磁気学の方程式が不完全なのだろう」と考え、光を伝える架空の媒質「エーテル」を探し求めました。しかしアインシュタインは逆転の発想をし、「電磁気学の教えの通り、真空中を進む光の速さは光源の運動に関係なく誰にとっても常に不変の $c$ である」と認め、むしろ200年間疑われなかったニュートンの「時間と空間の絶対性（ガリレイ変換）」のほうを書き直す道を選びました。この常識にとらわれない理性の勇気こそが、現代物理学の扉を開いたのです。
\end{examplebox}


\clearpage
\chapter{「同時性」の相対性}
\begin{center}
{\Large\bfseries\gtfamily 列車と堤防の思考実験が暴く絶対時間の崩壊}
\end{center}
\vspace*{2mm}

\section*{本章の視覚的見取り図}
\begin{center}
\begin{tikzpicture}[scale=0.9, >=Stealth]
    % Track and Light pulses
    \draw[thick, fill=gray!20] (-2,-0.2) rectangle (8,0.2);
    \node[draw, rectangle, fill=blue!10, thick, rounded corners, minimum width=8cm, minimum height=1.6cm] at (3,1.6) {};
    \node at (3, 2.2) {\textbf{高速列車（速度 $v \to$）}};
    % Light sources at A and B
    \fill[red] (-0.8, 0) circle (3pt) node[below=3pt] {落雷地点 $A$};
    \fill[red] (6.8, 0) circle (3pt) node[below=3pt] {落雷地点 $B$};
    \draw[decorate, decoration={snake, amplitude=2pt, segment length=5pt}, thick, orange!90!black, ->] (-0.7,0.4) -- (2.5,0.4) node[midway, above, font=\scriptsize] {光波 $c$};
    \draw[decorate, decoration={snake, amplitude=2pt, segment length=5pt}, thick, orange!90!black, ->] (6.7,0.4) -- (3.5,0.4) node[midway, above, font=\scriptsize] {光波 $c$};
    % Observers
    \node[draw, circle, fill=yellow!50, inner sep=2pt] at (3, 0) {$M$};
    \node[below=2pt, font=\scriptsize] at (3,-0.3) {堤防の中点 $M$（同時到達）};
    \node[draw, circle, fill=cyan!50, inner sep=2pt] at (3.8, 1.6) {$M'$};
    \node[above=2pt, font=\scriptsize] at (3.8, 2.0) {車内の中点 $M'$（$B$ の光に突進）};
    \draw[->, ultra thick, red] (4.2, 1.6) -- (5.0, 1.6) node[right, font=\scriptsize] {$v$};
\end{tikzpicture}

\small\textbf{図版解説：}列車と堤防の思考実験。堤防の中点 $M$ で同時の2つの落雷は、前進する車内の中点 $M'$ には同時に届かない。
\end{center}
\vspace*{2mm}

\section*{アインシュタインの原典思索}
\begin{origbox}[原著第3章：アルベルト・アインシュタインの肉声]
線路の堤防に沿って走る非常に長い列車を考えよう。堤防上の2地点 $A$ と $B$ に同時に落雷があったとする。堤防の中点 $M$ にいる観察者には、$A$ と $B$ からの光が同時に到達する。
しかし、列車のちょうど中央にいる乗客 $M'$ は、雷が落ちた瞬間には $M$ と一致していたとしても、光が伝わる間に列車とともに $B$ の方向へ進む。したがって $M'$ は $B$ からの光を迎えに行き、$A$ からの光から遠ざかる。
ゆえに $M'$ の目には、$B$ の閃光が $A$ の閃光よりも早く到達する。堤防の人にとって「同時」であった事象は、列車の乗客にとっては「同時ではない」のである。
\end{origbox}

\section*{大人のための徹底解説講義}
\begin{lecturebox}[徹底講義：高校数学で追体験する物理学革命]
物理学の歴史上、最も鮮やかで美しい思考実験が、この「列車と堤防の雷」です。
ニュートン力学では、「いま、この瞬間」という時間は、宇宙のどこであっても、誰にとっても寸分違わず等しく流れていると考えられていました。これを「絶対時間」と呼びます。

しかしアインシュタインは問います。「そもそも、遠く離れた2つの場所で起きた出来事が『同時』であるとは、具体的にどうやって確かめるのか？」
私たちが遠方の出来事を知るには、必ず「光」という信号を受け取る必要があります。そして光の速さは無限大ではなく、秒速30万キロメートルという有限の値です。
堤防に立っている人にとって同時に届いた雷の光は、前進している列車の乗客にとっては、自らが近づいていく前方（$B$）の光が先に届き、逃げていく後方（$A$）の光は遅れて届きます。光速が誰にとっても不変である以上、「前方の雷が先に落ちた」と結論づける以外にありません。
「同時」とは客観的・絶対的な事実ではなく、「誰が見ているか（どのような運動状態にあるか）」に完全に依存する相対的な概念だったのです。
\end{lecturebox}

\vspace*{3mm}
\begin{examplebox}{大人の学び直しのための実社会・重要例題：第3章}
\textbf{【実社会の具体例：国際金融取引・高頻度取引（HFT）におけるミリ秒単位の同時性】}\\
東京証券取引所、ニューヨーク証券取引所、ロンドン証券取引所を光海底ケーブルで結び、1秒間に数万件の自動売買を行う高頻度取引（HFT）では、「どちらの注文が先に取引所に届いたか」が数億円の損益を左右します。\\
(1) アインシュタインの列車と堤防の思考実験に基づき、「堤防の観測者にとって線路の両端 $A, B$ に同時に落ちた雷」が、列車内の乗客にはなぜ「前方の落雷 $B$ のほうが先に起きた」と観測されるのか、論理的理由を明快に説明せよ。\\
(2) なぜ「絶対的な同時性」という概念は錯覚であり、光速という有限の信号伝達速度が存在する宇宙においては「同時性は観測者の運動状態に依存する相対的なもの」にならざるを得ないのかを述べよ。
\\[2mm]
\textbf{【懇切丁寧な解答・ステップ解説】}\\
\textbf{(1) の解：}
堤防の中点にいる観測者 $M$ には、$A$ と $B$ から放たれた光が同時に届くため、「落雷は同時であった」と判定されます。
しかし、列車の中心にいる乗客 $M'$ は、雷が落ちた瞬間には $M$ と同じ位置にいたものの、光が進んでくる間に列車とともに前方（$B$ の方向）へ移動します。したがって、乗客 $M'$ は「前方の $B$ から来る光を迎えに行き、後方の $A$ から来る光からは逃げる」ことになります。その結果、光速度が不変である以上、$M'$ の目には明らかに前方の落雷 $B$ の閃光が先に飛び込んできます。すなわち、車内の人にとっては「$B$ が先に落ちた」のが客観的事実となります。
\\[2mm]
\textbf{(2) の解：}
もし情報や信号が一瞬で（無限大の速度で）宇宙の端から端まで届くなら、全宇宙共通の「神の絶対時間」を定義できます。しかし現実の宇宙には光速 $c$ という越えられない最高速度が存在します。「離れた2地点の出来事が同時である」と判定するには光信号による照合が不可欠であり、その観測者が動いていれば光の到達順序が入れ替わります。ゆえに、「絶対的な現在」などどこにも存在せず、同時性は観測者の視座に依存する相対的な概念なのです。
\end{examplebox}


\clearpage
\chapter{時間の遅れと「光時計」の数理}
\begin{center}
{\Large\bfseries\gtfamily 三平方の定理だけで導く動く時計の遅れの方程式}
\end{center}
\vspace*{2mm}

\section*{本章の視覚的見取り図}
\begin{center}
\begin{tikzpicture}[scale=0.95, >=Stealth]
    % Rest frame clock
    \begin{scope}[xshift=0cm]
        \draw[thick, fill=gray!10] (0,0) rectangle (2,3);
        \node at (1,3.3) {\textbf{静止系（固有時間 $\Delta t_0$）}};
        \fill[black] (1,0.3) circle (3pt) node[right, font=\scriptsize] {光源/検出器};
        \draw[ultra thick, blue!80!black] (0.3,2.7) -- (1.7,2.7) node[midway, above, font=\scriptsize] {鏡};
        \draw[decorate, decoration={snake, amplitude=2pt, segment length=6pt}, thick, red, ->] (0.9,0.4) -- (0.9,2.6);
        \draw[decorate, decoration={snake, amplitude=2pt, segment length=6pt}, thick, red, ->] (1.1,2.6) -- (1.1,0.4);
        \draw[<->, thick] (0.1,0.3) -- (0.1,2.7) node[midway, left] {$L$};
        \node[below=4pt, font=\footnotesize] at (1,0) {$\Delta t_0 = \dfrac{2L}{c}$};
    \end{scope}
    % Moving frame clock
    \begin{scope}[xshift=4.5cm]
        \draw[->, thick, gray!60] (-0.5,0) -- (5,0) node[right, black] {$x$};
        \node at (2,3.3) {\textbf{移動系（光が斜めに進む $\Delta t$）}};
        \fill[black] (0,0.3) circle (2.5pt);
        \draw[ultra thick, blue!80!black] (1.2,2.7) -- (2.6,2.7);
        \fill[black] (3.8,0.3) circle (2.5pt);
        \draw[dashed, red, thick] (0,0.3) -- (1.9,2.7) -- (3.8,0.3);
        \draw[->, thick, red] (0,0.3) -- (0.95,1.5);
        \draw[->, thick, red] (1.9,2.7) -- (2.85,1.5);
        \draw[<->, blue, thick] (0,-0.4) -- (3.8,-0.4) node[midway, below, font=\scriptsize] {$v \Delta t$};
        \node[align=left, fill=white, draw=gray!40, rounded corners, font=\scriptsize] at (2, 1.2) {三平方の定理：\\$\left(\frac{c\Delta t}{2}\right)^2 = L^2 + \left(\frac{v\Delta t}{2}\right)^2$};
    \end{scope}
\end{tikzpicture}

\small\textbf{図版解説：}光時計の思考実験。三平方の定理から時間の遅れ $\Delta t = \Delta t_0 / \sqrt{1 - v^2/c^2}$ が導かれる幾何学的構造。
\end{center}
\vspace*{2mm}

\section*{アインシュタインの原典思索}
\begin{origbox}[原著第4章：アルベルト・アインシュタインの肉声]
光速度がすべての慣性系で一定であるという公理から、運動する系における時間の進み方を直接導き出すことができる。
運動している座標系 $K'$ の中に置かれた時計は、その運動を外の静止座標系 $K$ から観測したとき、静止している同一の時計に比べて遅れて進む。この遅れの比率は、速度比 $v/c$ を用いて $\sqrt{1 - v^2/c^2}$ という因子によって正確に決定される。
\end{origbox}

\section*{大人のための徹底解説講義}
\begin{lecturebox}[徹底講義：高校数学で追体験する物理学革命]
「動いているものの時間は遅れる」という有名な相対論の結論は、実は高校1年生で習う「三平方の定理（ピタゴラスの定理）」だけで完全に証明できます。
床と天井に鏡を向かい合わせに置き、その間を光の粒が上下に往復する「光時計」を想像してください。
車内に座っている乗客から見れば、光は天井に向かって真っ直ぐ垂直に往復するだけです（往復距離 $2L$）。

しかし、駅のホームに立っている人からこの列車を見ると、光が天井に到達するまでの間に列車が横へ距離 $v \Delta t$ だけ走るため、光の軌跡は「斜めの二等辺三角形の山」を描きます。
三角形の斜辺は、垂直の高さ $L$ よりも明らかに長くなります。そして、ホームの人にとっても光の速度は車内と同じ $c$ でなければなりません。「より長い距離を、同じ速さで進む」のですから、ホームの人から見れば、光が往復するのに「より長い時間（$\Delta t$）」がかかることになります。
これを数式で整理すると、
\[
\Delta t = \frac{\Delta t_0}{\sqrt{1 - \frac{v^2}{c^2}}}
\]
という美しい公式が導かれます。速度 $v$ が光速 $c$ に近づくほど、分母はゼロに近づき、時間の遅れは無限大へと拡大していきます。
\end{lecturebox}

\vspace*{3mm}
\begin{examplebox}{大人の学び直しのための実社会・重要例題：第4章}
\textbf{【実社会の具体例：航空機に原子時計を載せて世界一周したハフェル・キーティング実験】}\\
1971年、物理学者ジョセフ・ハフェルとリチャード・キーティングは、セシウム原子時計を定期航空機に乗せて地球を東西方向に世界一周飛行させ、地上の静止時計との時間の進み方を比較しました。\\
(1) 光時計の思考実験において、三平方の定理から導かれた「動く時計の遅れ $\Delta t = \Delta t_0 / \sqrt{1 - v^2/c^2}$」の物理的意味を、光の軌跡が斜めに伸びる幾何学的関係から解説せよ。\\
(2) ハフェル・キーティング実験において、地球の自転に伴う東回りと西回りの速度差によって、実際にナノ秒単位で特殊相対論の遅れが実測された意義を述べよ。
\\[2mm]
\textbf{【懇切丁寧な解答・ステップ解説】}\\
\textbf{(1) の解：}
列車の中に垂直に立てた鏡の間を光が往復する時計（光時計）を考えます。車内の人にとっては光は垂直に距離 $2L$ を往復するだけです。
しかし、外にいる人から見ると、光が上下に移動する間に列車が水平に走るため、光の軌道は「V字形の斜めの軌跡」を描きます。斜辺は垂直の高さより必ず長くなります。外の人にとっても光の進む速さは同じ $c$ ですから、より長い距離を進むためには「より多くの時間」がかからねばなりません。つまり、車内での $1$ 秒の往復が、外の人から見れば $1$ 秒以上の時間がかかって見えることになり、「動いている乗り物の時計は外から見るとゆっくり進む」という結論が幾何学的に確定します。
\\[2mm]
\textbf{(2) の解：}
この実験により、旅客機という私たちの身近な乗り物の速度（時速約 $900\,\text{km}$）であっても、アインシュタインの時間の遅れが机上の空論ではなく、物理的現実として時計の刻みを遅らせていることが完全に証明されました。時間は単なる主観的な感覚ではなく、空間の運動と結びついた客観的な物理量であることが実証された歴史的転換点です。
\end{examplebox}


\clearpage
\chapter{ローレンツ収縮と動く物体の縮み}
\begin{center}
{\Large\bfseries\gtfamily 進行方向の空間が物理的に収縮するメカニズム}
\end{center}
\vspace*{2mm}

\section*{本章の視覚的見取り図}
\begin{center}
\begin{tikzpicture}[scale=0.9, >=Stealth]
    % Rest rod
    \node[anchor=east, font=\bfseries] at (-0.5,2.5) {静止時 $L_0$：};
    \draw[fill=blue!20, draw=blue!80, thick] (0,2.1) rectangle (6,2.9);
    \draw[|<->|, thick] (0,3.2) -- (6,3.2) node[midway, above] {固有長 $L_0$};
    % Moving rod
    \node[anchor=east, font=\bfseries] at (-0.5,0.7) {高速運動時 $L$：};
    \draw[fill=red!20, draw=red!80, thick] (1,0.3) rectangle (4.6,1.1);
    \draw[|<->|, thick] (1,-0.2) -- (4.6,-0.2) node[midway, below] {$L = L_0 \sqrt{1 - v^2/c^2}$};
    \draw[->, ultra thick, red!80!black] (4.7,0.7) -- (6.2,0.7) node[right] {速度 $v$};
    % Contraction note
    \node[align=center, fill=yellow!10, draw=orange!60, rounded corners, font=\footnotesize] at (7.5, 2.0) {ローレンツ収縮：\\運動方向のみ縮み、\\垂直方向の幅は不変};
\end{tikzpicture}

\small\textbf{図版解説：}進行方向に沿って剛体の長さが縮むローレンツ収縮の幾何学的イメージ。
\end{center}
\vspace*{2mm}

\section*{アインシュタインの原典思索}
\begin{origbox}[原著第5章：アルベルト・アインシュタインの肉声]
静止しているときに長さ $L_0$ を持つ剛体の定規を考える。この定規が自身の長さの方向に沿って速度 $v$ で運動しているとき、静止系 $K$ からその長さを測定すると、定規の長さは
\[
L = L_0 \sqrt{1 - \frac{v^2}{c^2}}
\]
に短縮されて観測される。運動が速ければ速いほど物体は運動方向に短くなり、仮に光速 $c$ に達すれば厚みはゼロとなってしまう。ただし進行方向に垂直な幅は全く変化しない。
\end{origbox}

\section*{大人のための徹底解説講義}
\begin{lecturebox}[徹底講義：高校数学で追体験する物理学革命]
時間が伸びるのと表裏一体の関係にあるのが、「空間（長さ）の縮み」です。これをローレンツ収縮と呼びます。
時速 $300\,\text{km}$ の新幹線であっても、理論上は静止しているときよりわずかに縮んでいます（日常の速度では $10^{-15}$ 程度という原子核サイズの微小量ですが、加速器の中の素粒子では劇的に縮みます）。

なぜ物体が縮むのでしょうか？ 物体が押しつぶされるような力が働いているわけではありません。
「動いている物体の長さを測る」とは、「物体の前端と後端の位置を【同時に】記録してその差をとる」という測定行為です。ところが第3章で見たように、「同時」の基準そのものが運動状態によってズレてしまいます。
時間の基準がズレる結果として、静止系から測定した物体の長さは、必ず $\sqrt{1 - v^2/c^2}$ 倍に縮んで測定されるのです。
空間とは硬い不変のステージではなく、時間をどのように測るかによって伸縮する柔軟な布地のようなものなのです。
\end{lecturebox}

\vspace*{3mm}
\begin{examplebox}{大人の学び直しのための実社会・重要例題：第5章}
\textbf{【実社会の具体例：高エネルギー加速器での素粒子ビームとローレンツ収縮】}\\
スイス・ジュネーブのCERN（欧州原子核研究機構）にある大型ハドロン衝突型加速器（LHC）では、陽子ビームを光速の $99.9999991\%$ まで加速して正面衝突させます。\\
(1) 静止した観測者から見て、光速に近い速度で飛ぶ陽子の球体が進行方向に沿ってどのように変形して見えるか、ローレンツ収縮の式 $L = L_0 \sqrt{1 - v^2/c^2}$ を用いて説明せよ。\\
(2) なぜローレンツ収縮は物体の横幅（進行方向に垂直な方向）には生じず、進行方向の長さだけに起きるのか、座標変換の直交性から考察せよ。
\\[2mm]
\textbf{【懇切丁寧な解答・ステップ解説】}\\
\textbf{(1) の解：}
静止状態では陽子は真球（球形）をしています。しかし光速の $99.9999991\%$（ローレンツ因子 $\gamma \approx 7450$）に達すると、実験室の観測者から見た陽子の進行方向の厚みは、静止時の約 $\frac{1}{7450}$ にまで激しく押し縮められます。すなわち、球体ではなく「極めて薄い平たいパンケーキ（円盤状）」のような形状になって衝突点に突入します。加速器の検出器設計では、このローレンツ収縮による電荷密度の増大を考慮しなければ正確な衝突解析ができません。
\\[2mm]
\textbf{(2) の解：}
運動は $x$ 軸方向（進行方向）にのみ起きており、$y$ 軸や $z$ 軸（進行方向に垂直な方向）には相対運動がありません。もし進行方向に垂直な方向の長さも変わってしまうと、すれ違う2本の同じ長さの棒について「互いに相手のほうが細い」と主張し合い、すり抜けられるかどうかの物理的事実が破綻してしまいます。したがって、収縮は相対運動が起きている進行軸のみに厳密に限定されます。
\end{examplebox}


\clearpage
\chapter{質量とエネルギーの等価性：$E=mc^2$}
\begin{center}
{\Large\bfseries\gtfamily 光の圧力が明かす質量欠損と核エネルギーの源泉}
\end{center}
\vspace*{2mm}

\section*{本章の視覚的見取り図}
\begin{center}
\begin{tikzpicture}[scale=0.9, >=Stealth]
    % Mass increase curve
    \draw[->, thick] (0,0) -- (6.5,0) node[right] {速度 $v$};
    \draw[->, thick] (0,0) -- (0,4.2) node[above] {相対論的質量 $m(v)$ / エネルギー $E$};
    \draw[dashed, red, thick] (5.5,0) -- (5.5,4.0) node[above, font=\footnotesize] {光速壁 $c$};
    \draw[domain=0:5.2, smooth, variable=\x, ultra thick, blue!80!black] plot ({\x}, {1.0 / sqrt(1 - (\x/5.5)*(\x/5.5))});
    \fill[black] (0,1.0) circle (2.5pt) node[left, font=\footnotesize] {静止質量 $m_0$ ($E_0 = m_0 c^2$)};
    \node[align=left, fill=white, draw=gray!40, rounded corners, font=\footnotesize] at (2.8, 3.2) {速度が光速に近づくと\\質量・エネルギーは無限大に発散：\\物体は光速を超えられない};
\end{tikzpicture}

\small\textbf{図版解説：}質量とエネルギーの等価性 $E = mc^2$ と、光速 $c$ へ近づくにつれて無限大へ発散する相対論的質量。
\end{center}
\vspace*{2mm}

\section*{アインシュタインの原典思索}
\begin{origbox}[原著第6章：アルベルト・アインシュタインの肉声]
物体が放射の形でエネルギー $L$ を放出するとき、その物体の質量は $L/c^2$ だけ減少する。不活性な質量とエネルギーとの間には本質的な差異は存在しない。
物体のエネルギーが変化すれば、それに伴って物体の質量も変化する。静止質量 $m$ を持つ物体は、その内部に
\[
E = mc^2
\]
に等しい莫大なエネルギーを蓄えている。放射性物質の崩壊において観測される莫大な熱の発生は、この質量欠損がエネルギーへと姿を変えた直接の証拠である。
\end{origbox}

\section*{大人のための徹底解説講義}
\begin{lecturebox}[徹底講義：高校数学で追体験する物理学革命]
世界で最も有名な物理方程式、$E = mc^2$。この簡潔を極めた式の真の意味を、アインシュタインは1905年の補遺論文で論証しました。
物体が光（電磁波）を前後に放出したとき、エネルギー保存則と運動量保存則を満たすためには、光を放出したあとの物体の「慣性（動かしにくさ）」がわずかに軽くなっていなければならないことを突き止めたのです。

光速 $c$ は秒速30万キロメートル、すなわち $3 \times 10^8\,\text{m/s}$ です。その2乗である $c^2$ は
\[
c^2 \approx 9 \times 10^{16} \approx 90,000,000,000,000,000
\]
という天文学的・圧倒的な巨大さです。これは、わずか $1\,\text{グラム}$ の物質が完全にエネルギーに変換されただけで、$90$ 兆ジュール（大型石油タンカー数隻分の燃料を燃やした熱量に匹敵）のエネルギーが解き放たれることを意味します。
質量とは物質の「重さ」という固定的な属性ではなく、「超高密度に凝縮されたエネルギーそのものの姿」だったのです。
\end{lecturebox}

\vspace*{3mm}
\begin{examplebox}{大人の学び直しのための実社会・重要例題：第6章}
\textbf{【実社会の具体例：原子力発電所・PETがん検診と質量・エネルギー等価性 $E=mc^2$】}\\
現在、日本の先端医療施設で広く行われているPET検査（陽電子放射断層撮影）では、体内に注入した薬剤から放出される「陽電子」が体内の「電子」と衝突し、両者が跡形もなく消滅して2本のガンマ線（光子エネルギー）に変わる現象を利用しています。\\
(1) 質量とエネルギーの等価式 $E = mc^2$ の意味を、電子と陽電子の対消滅（質量が $100\%$ 光エネルギーに変換される現象）を例に説明せよ。\\
(2) なぜ古典物理学では「質量保存の法則」と「エネルギー保存の法則」が別々の法則として信じられていたのか、石炭の燃焼などの化学反応で放出されるエネルギーの微小さの観点から考察せよ。
\\[2mm]
\textbf{【懇切丁寧な解答・ステップ解説】}\\
\textbf{(1) の解：}
$E = mc^2$ は、「質量とは、閉じ込められた莫大なエネルギーの形態の一つにすぎない」という宇宙の根本真理を示しています。PET検査では、電子（質量 $m$）と陽電子（質量 $m$）が出会うと、物質としての質量が完全に消滅し、合計 $2mc^2 \approx 1.022\,\text{MeV}$ の高エネルギー電磁波（ガンマ線）となって正反対の方向に放射されます。PET装置はこの光子をリング状検出器で捉えてがん細胞の位置をミリ単位で特定します。医療現場で日常的に $E=mc^2$ の直訳的応用が使われているのです。
\\[2mm]
\textbf{(2) の解：}
石炭やガソリンを燃やす化学反応でも、厳密には発生した熱エネルギーの分だけ質量が減少（質量欠損）しています。しかし、$c^2 \approx 9 \times 10^{16}$ という変換係数が途方もなく巨大であるため、放出される熱に対する質量減少は全体の $100億分の1$ 以下という、化学天秤では到底検知できない極微量にすぎませんでした。そのため、19世紀までの科学者は質量とエネルギーが独立に保存していると誤認していたのです。
\end{examplebox}


\clearpage
\chapter{ミンコフスキーの4次元時空連続体}
\begin{center}
{\Large\bfseries\gtfamily 時間と空間の統合と「世界線」の幾何学}
\end{center}
\vspace*{2mm}

\section*{本章の視覚的見取り図}
\begin{center}
\begin{tikzpicture}[scale=1.0, >=Stealth]
    % Coordinate axes
    \draw[->, thick] (-3,0) -- (3,0) node[right] {$x$ (空間)};
    \draw[->, thick] (0,-0.5) -- (0,3.5) node[above] {$ct$ (時間)};
    % Light cone
    \fill[yellow!20, opacity=0.7] (0,0) -- (3,3) -- (-3,3) -- cycle;
    \draw[ultra thick, orange!90!black] (-3,3) -- (0,0) -- (3,3);
    \node[above, orange!90!black, font=\footnotesize] at (2.4, 2.5) {光世界線 ($x = ct$)};
    \node[above, orange!90!black, font=\footnotesize] at (-2.4, 2.5) {光世界線 ($x = -ct$)};
    % Worldline of particle
    \draw[ultra thick, blue!80!black] (0,0) .. controls (0.5,1.2) and (-0.3,2.0) .. (0.8,3.2) node[above, font=\footnotesize] {物質の世界線 ($v < c$)};
    % Spacetime invariant
    \node[align=center, fill=white, draw=gray!50, rounded corners, font=\footnotesize] at (0, 1.8) {絶対的不変間隔：\\$s^2 = c^2 t^2 - x^2 > 0$\\（未来光円錐の内部）};
\end{tikzpicture}

\small\textbf{図版解説：}ミンコフスキーの4次元時空と光円錐（Light Cone）。すべての因果律は光円錐の内部に制限される。
\end{center}
\vspace*{2mm}

\section*{アインシュタインの原典思索}
\begin{origbox}[原著第7章：アルベルト・アインシュタインの肉声]
数学者ヘルマン・ミンコフスキーは1908年にこう宣言した：
「今日以後、空間それ自身、時間それ自身は完全に影の中に消え去り、両者の或る種の結合のみが独立した実在を保ち続けるであろう」。
3次元空間の座標 $(x, y, z)$ に虚数時間 $ict$（あるいは時間軸 $ct$）を加えた4次元連続体において、ローレンツ変換は4次元空間における座標軸の直交回転として極めて自然に理解される。
\end{origbox}

\section*{大人のための徹底解説講義}
\begin{lecturebox}[徹底講義：高校数学で追体験する物理学革命]
アインシュタインのかつての数学の恩師であったヘルマン・ミンコフスキーは、特殊相対性理論の数理を一歩進め、現代物理学の基本言語となる「4次元時空（Spacetime）」の概念を創始しました。
3次元空間の中で、私たちが斜めからサイコロを眺めると、見る角度によって「横幅」が広く見えたり「奥行き」が狭く見えたりします。しかし、サイコロそのものの立体的な大きさは変わっていません。

ミンコフスキーは、相対論における「時間の遅れ」や「物体の収縮」もこれと全く同じ現象であると喝破しました。
動いている観測者とは、4次元時空の中で「座標軸が斜めに傾いている観測者」にすぎません。傾いた軸で測るから、時間成分が伸びて見えたり、空間成分が縮んで見えたりするのです。
空間の隔たり $\Delta x$ と時間の隔たり $\Delta t$ を組み合わせた不変量
\[
s^2 = c^2 (\Delta t)^2 - (\Delta x)^2 - (\Delta y)^2 - (\Delta z)^2
\]
は、いかなる等速運動をする観測者から見ても完全に一致します。時間と空間は別個のものではなく、「4次元時空連続体」という織物の縦糸と横糸だったのです。
\end{lecturebox}

\vspace*{3mm}
\begin{examplebox}{大人の学び直しのための実社会・重要例題：第7章}
\textbf{【実社会の具体例：光円錐と宇宙の因果律（過去・現在・未来の境界線）】}\\
天文学者が「130億光年彼方の初期銀河の光」をジェームズ・ウェッブ宇宙望遠鏡で観測するとき、私たちは「130億年前の過去の姿」を今この瞬間に見ています。\\
(1) ミンコフスキーの4次元時空において、光円錐（過去光円錐・未来光円錐）がなぜ物理学における「因果関係（原因と結果）」の絶対の防壁となっているのかを説明せよ。\\
(2) 光速を超える情報伝達（超光速通信）が存在しない限り、なぜ「過去に戻って歴史を変えてしまうタイムパラドックス」が宇宙において厳格に禁止されているのかを考察せよ。
\\[2mm]
\textbf{【懇切丁寧な解答・ステップ解説】}\\
\textbf{(1) の解：}
あらゆる物質や信号は光速 $c$ を超えて移動できません。ある事象（出来事）から発した影響が届く範囲は、時間とともに光速で広がる「未来光円錐の内部」に限られます。同様に、現在の私たちに影響を及ぼし得た過去の事象はすべて「過去光円錐の内部」に収まっています。光円錐の外側（他所：Elsewhere）にある事象とは、光速をもってしても情報が届かないため、物理的に原因にも結果にもなり得ません。ミンコフスキー時空は因果の連鎖を幾何学的に整然と隔離しているのです。
\\[2mm]
\textbf{(2) の解：}
もし光速を超える信号が存在すると、ローレンツ変換によって「ある観測者にとっては原因より先に結果が起きる（過去への情報到達）」という逆転現象が生じてしまいます。これは親が生まれる前に子どもが親を殺害できてしまうような因果律の破綻を招きます。宇宙が自己矛盾を起こさず論理的一貫性を保ち続けているのは、光速度という絶対の速度制限が存在するおかげなのです。
\end{examplebox}


\clearpage
\part{一般相対性理論（重力と曲がった時空）}

\chapter{一般相対性理論への道と「等価原理」}
\begin{center}
{\Large\bfseries\gtfamily エレベーターの思考実験と重力＝加速度の発見}
\end{center}
\vspace*{2mm}

\section*{本章の視覚的見取り図}
\begin{center}
\begin{tikzpicture}[scale=0.9, >=Stealth]
    % Earth gravitational field
    \begin{scope}[xshift=0cm]
        \draw[thick, fill=blue!10] (0,0) rectangle (2.5,3.2);
        \node at (1.25,3.5) {\textbf{重力場（地上）}};
        \fill[gray!40] (-0.3,-0.2) rectangle (2.8,0);
        \node at (1.25,-0.4) {地球表面};
        \fill[red] (1.25,1.8) circle (3pt) node[right, font=\footnotesize] {自由落下する球};
        \draw[->, ultra thick, red] (1.25,1.7) -- (1.25,0.7) node[midway, right, font=\footnotesize] {$g$};
        \node[font=\footnotesize] at (1.25, 0.4) {重力加速度};
    \end{scope}
    % Equal sign
    \node[font=\Huge\bfseries] at (3.5, 1.8) {$\equiv$};
    % Accelerating rocket
    \begin{scope}[xshift=4.5cm]
        \draw[thick, fill=yellow!10] (0,0) rectangle (2.5,3.2);
        \node at (1.25,3.5) {\textbf{無重力空間の加速ロケット}};
        \draw[->, ultra thick, purple] (1.25,3.2) -- (1.25,4.2) node[above, font=\footnotesize] {上向き加速 $\alpha = g$};
        \fill[red] (1.25,1.8) circle (3pt) node[right, font=\footnotesize] {静止する球};
        \draw[->, ultra thick, red] (1.25,1.7) -- (1.25,0.7) node[midway, right, font=\footnotesize] {見かけの落下 $g$};
        \node[font=\footnotesize] at (1.25, 0.4) {慣性力（擬似重力）};
    \end{scope}
\end{tikzpicture}

\small\textbf{図版解説：}アインシュタインの等価原理。一様な重力場と等加速度運動する乗り物の中の物理現象は完全に等価である。
\end{center}
\vspace*{2mm}

\section*{アインシュタインの原典思索}
\begin{origbox}[原著第8章：アルベルト・アインシュタインの肉声]
広大な宇宙空間に浮かぶ密閉された大きな箱（エレベーター）を想定しよう。その天井の外側にロープがつけられ、何者かが一定の力で箱を引き上げているとする。箱の中の人間は、床に向かって放した物体がすべて等しい加速度で落下するのを目撃する。
箱の中の人は「自分たちは地上の一様な重力場の中に静止している」と信じ込むであろう。そしていかなる実験を行っても、それが「外からの加速による慣性力」なのか「真の重力」なのかを区別することは原理的に不可能である。これが等価原理である。
\end{origbox}

\section*{大人のための徹底解説講義}
\begin{lecturebox}[徹底講義：高校数学で追体験する物理学革命]
特殊相対性理論は「等速直線運動をしている観測者（慣性系）」だけを扱っていました。しかし、現実の宇宙には加速する乗り物もあれば、天体が及ぼし合う重力も満ち満ちています。
「あらゆる運動、加速運動も含めて物理法則が同じになるような【一般相対性理論】は創れないだろうか？」
アインシュタインはこの難問に10年近く挑み続けました。その突破口となったのが、1907年に閃いた「等価原理」です。

ビルのエレベーターのワイヤーが切れて自由落下したとき、中にいる人はフワリと宙に浮き、体重計の針はゼロを指します（無重力）。逆に、宇宙空間でロケットが $9.8\,\text{m/s}^2$ で急加速すると、乗組員は地球と同じ重力で床に押しつけられます。
アインシュタインは直観しました。「重力（引力）と、加速によって生じる力（慣性力）は、本質的に全く同じものである！」
この等価原理を出発点とすることで、アインシュタインは「重力とは何か？」というニュートン以来の最大の謎に挑むことになります。
\end{lecturebox}

\vspace*{3mm}
\begin{examplebox}{大人の学び直しのための実社会・重要例題：第8章}
\textbf{【実社会の具体例：ビルのエレベーターのワイヤーが切れたら？（無重力の創出）】}\\
国際宇宙ステーション（ISS）の中の宇宙飛行士が宙に浮いているのは、宇宙空間に地球の重力がないからではなく、ISS全体が秒速 $7.7\,\text{km}$ で地球の周りを「自由落下し続けている」からです。\\
(1) アインシュタインが生涯で「最も幸福な考え（der glücklichste Gedanke meines Lebens）」と呼んだ、「屋根から落下している人には自らの体重が感じられない」という着想が、どのようにして一般相対論の等価原理へと結実したのかを述べよ。\\
(2) 局所的な自由落下系の中では、重力という「力」が完全に消去され、特殊相対性理論の慣性系と全く区別がつかなくなる論理を説明せよ。
\\[2mm]
\textbf{【懇切丁寧な解答・ステップ解説】}\\
\textbf{(1) の解：}
1907年、スイス特許局で執務していたアインシュタインは、窓の外を見ながら「もし人がビルの屋上から自由落下したら、その人は自分の体重を感じない（無重力になる）」と気づきました。落下するエレベーターの中では、手から離したリンゴは床に落ちず、空中に静止して見えます。重力によって引き落とされる力と、落下運動に伴う慣性力が完全に相殺するからです。アインシュタインはこの直観から、「重力とは外的な引っ張る力ではなく、加速系における慣性力と全く同一の現象である」という等価原理を打ち立てました。
\\[2mm]
\textbf{(2) の解：}
箱の外が見えない密閉された実験室の中にいる限り、いかなる精密な物理実験（振り子、物体の落下、天秤）を行っても、「自分たちが深宇宙の無重力空間で静止しているのか」、それとも「重力場の中で箱ごと自由落下しているのか」を判定することは原理的に不可能です。すなわち、局所的には重力場は座標系の取り方（自由落下）によって消し去ることができるため、重力を幾何学的な空間の歪みとして再定式化する道が拓かれたのです。
\end{examplebox}


\clearpage
\chapter{回転円盤と非ユークリッド幾何学の必然性}
\begin{center}
{\Large\bfseries\gtfamily 円周率が3.14ではなくなる世界と曲がった時空}
\end{center}
\vspace*{2mm}

\section*{本章の視覚的見取り図}
\begin{center}
\begin{tikzpicture}[scale=0.95, >=Stealth]
    % Rotating disk
    \draw[thick, fill=cyan!10] (0,0) circle (2.5);
    \draw[->, ultra thick, red] (1.8, 1.8) arc (45:85:2.5) node[above, font=\footnotesize] {角速度 $\omega$};
    \fill[black] (0,0) circle (2pt) node[below=2pt, font=\footnotesize] {中心 $O$（静止）};
    % Radial rod (no contraction)
    \draw[|-|, thick, blue] (0,0) -- (0,2.3) node[midway, left, font=\scriptsize] {半径 $R$（収縮なし）};
    % Circumference rods (Lorentz contracted)
    \draw[ultra thick, red!80!black] (2.3, -0.4) arc (-10:30:2.3) node[midway, right, font=\scriptsize] {周長方向（速度 $v = R\omega$ で収縮）};
    % Non-euclidean conclusion
    \node[align=left, fill=white, draw=gray!50, rounded corners, font=\footnotesize] at (0, -3.2) {円周上の定規のみがローレンツ収縮するため：\\$\dfrac{\text{測定円周}}{\text{測定直径}} > \pi$ （ユークリッド幾何学の崩壊！）};
\end{tikzpicture}

\small\textbf{図版解説：}回転円盤の思考実験。円周方向の定規のみが収縮するため、円周率が $\pi$ より大きくなり非ユークリッド幾何学が現れる。
\end{center}
\vspace*{2mm}

\section*{アインシュタインの原典思索}
\begin{origbox}[原著第9章：アルベルト・アインシュタインの肉声]
中心のまわりを高速回転する円盤を考えよう。円盤の半径方向に沿って置かれた定規は運動方向に垂直であるから収縮しない。しかし円盤の円周に沿って置かれた定規は、接線方向の運動に伴ってローレンツ収縮を起こす。
したがって、円盤の縁に沿って定規を敷き詰めて円周を実測すると、静止時の計算値 $2\pi R$ よりも多くの定規が必要となる。すなわち測定される円周と直径の比は $\pi$ よりも大きくなってしまう。回転系（加速度系、すなわち等価原理による重力場）においては、ユークリッド幾何学はもはや厳密には成立しない。
\end{origbox}

\section*{大人のための徹底解説講義}
\begin{lecturebox}[徹底講義：高校数学で追体験する物理学革命]
幾何学の常識を覆す極めつけの思考実験が「回転する円盤」です。
平らな円盤を高速で回転させ、その円周と直径を小さな定規を並べて測ることを考えます。
中心から縁に向かう「半径」を測るとき、定規の向きは回転の進行方向に対して直角ですから、定規は縮みません。
ところが、ぐるりと「円周」を測るとき、定規は進行方向に沿って並べられるため、ローレンツ収縮を起こして短くなってしまいます。

短くなった定規で円周を測れば、定規の目盛りの数は当然多くなります。その結果、
\[
\frac{\text{測定された円周}}{\text{測定された直径}} > 3.14159\dots (\pi)
\]
という、古代ギリシャのユークリッド幾何学では絶対にあり得ない事態が発生します。
回転運動（加速度）が存在する場所では、空間そのものが歪んでおり、曲面を扱う「非ユークリッド幾何学（リーマン幾何学）」を使わなければならないのです。
等価原理により「加速度＝重力」ですから、この結論は直ちに「重力がある場所の時空は曲がっている」という驚くべき真実を指し示していました。
\end{lecturebox}

\vspace*{3mm}
\begin{examplebox}{大人の学び直しのための実社会・重要例題：第9章}
\textbf{【実社会の具体例：遊園地の回転コーヒーカップと円周率 $\pi$ の破綻】}\\
遊園地の高速回転する大型アトラクションに乗ると、強い遠心力（外向きの力）を感じます。\\
(1) アインシュタインの回転円盤の思考実験において、円盤の縁に沿って置かれた定規がローレンツ収縮する一方、中心から外へ向かう半径方向の定規は収縮しない理由を説明せよ。\\
(2) その結果、測定された円周 $C$ と直径 $D$ の比 $C/D$ が $3.14159\dots$（$\pi$）よりも大きくなってしまうという幾何学的パラドックスから、なぜ「加速度や重力が存在する世界では非ユークリッド幾何学を採用せざるを得ない」と結論づけられるのかを考察せよ。
\\[2mm]
\textbf{【懇切丁寧な解答・ステップ解説】}\\
\textbf{(1) の解：}
角速度 $\omega$ で回転する円盤を外の静止系から眺めます。半径方向の定規は運動方向に対して垂直であるため、ローレンツ収縮せず、静止時と同じ長さ $R$ を保ちます。
しかし、円周に沿って並べられた定規は、接線方向の速度 $v = R\omega$ で高速運動しているため、運動方向に沿ってローレンツ収縮を起こし、目盛りの長さが短くなります。
\\[2mm]
\textbf{(2) の解：}
短くなった定規を使って円周を測ると、定規が縮んでいる分だけ、より多くの本数が必要になります。したがって、測定される円周の値 $C$ は、静止時の幾何学的計算値 $2\pi R$ よりも必然的に大きくなります。その結果、
\[
\frac{C}{2R} > \pi
\]
となり、古代ギリシャ以来信じられてきたユークリッド幾何学の公理「円周と直径の比は常に $\pi$ である」が完全に破綻します。等価原理によれば、回転に伴う遠心力（加速度）は重力と等価です。ゆえに、重力が作用している現実の空間は、平坦なユークリッド空間ではなく、曲がった「非ユークリッド幾何学（リーマン幾何学）」で記述されねばならないという歴史的結論に達したのです。
\end{examplebox}


\clearpage
\chapter{一般相対論の実験的検証}
\begin{center}
{\Large\bfseries\gtfamily 光の重力湾曲・水星の近日点移動・時間の遅れの実証}
\end{center}
\vspace*{2mm}

\section*{本章の視覚的見取り図}
\begin{center}
\begin{tikzpicture}[scale=0.95, >=Stealth]
    % Massive star / Sun
    \shade[inner color=yellow!80!white, outer color=orange!90!black] (2,1) circle (0.9);
    \node at (2,1) {\textbf{太陽 $M$}};
    % Distant star real position
    \fill[blue] (-2,3) circle (2.5pt) node[above left, font=\footnotesize] {恒星の真の位置};
    % Deflected light path
    \draw[thick, blue, dashed] (-2,3) -- (5.5,1.9) node[right, font=\footnotesize] {見かけの位置};
    \draw[ultra thick, red, ->] (-2,3) .. controls (1.5,1.9) and (2.5,1.9) .. (5,0.5) node[right, black, font=\footnotesize] {地球の観測者};
    % Deflection angle theta
    \draw[<->, thick, purple] (4.5,1.0) arc (-10:20:1.5) node[midway, right, font=\footnotesize] {偏向角 $\theta = \dfrac{4GM}{c^2 R}$};
\end{tikzpicture}

\small\textbf{図版解説：}重力場による光線の湾曲（重力レンズ効果）。太陽の質量が周囲の時空を歪め、背後の恒星光が曲げられて地球に届く。
\end{center}
\vspace*{2mm}

\section*{アインシュタインの原典思索}
\begin{origbox}[原著第10章：アルベルト・アインシュタインの肉声]
一般相対性理論は3つの決定的な観測的検証を提示した。
第1に、太陽の近傍を通過する恒星の光線が、太陽の重力場によって外側へ $1.75$ 秒角偏向すること。これは1919年の皆既日食観測によって見事に確認された。
第2に、ニュートン力学では説明できなかった水星の近日点が、1世紀あたり $43$ 秒角余分に回転することの完全な理論的説明。
第3に、強い重力場から放射される光の波長が長くなる重力赤方偏移（時計の進みの遅れ）である。
\end{origbox}

\section*{大人のための徹底解説講義}
\begin{lecturebox}[徹底講義：高校数学で追体験する物理学革命]
どれほど美しい理論であっても、実験と観測によって検証されなければ物理学の真理とは認められません。アインシュタインは自らの理論の命運を賭けて、3つの具体的な予言を行いました。

その筆頭が「太陽による光の曲がり」です。一般相対論によれば、太陽の巨大な質量によって周囲の時空がすり鉢状に凹んでいるため、その縁をかすめる星の光はレンズのように曲げられます。
1919年、イギリスの天文学者アーサー・エディントン率いる観測隊が西アフリカとブラジルで皆既日食を観測し、太陽の背後にある恒星の位置が、アインシュタインの予言通り正確に $1.75$ 秒角ズレていることを撮影乾板から証明しました。

さらに、天文学者を100年間悩ませていた「水星の軌道が少しずつ回ってしまう謎（近日点移動）」についても、アインシュタイン方程式は見事に余剰の43秒角を完璧に弾き出しました。
ニュートン力学の限界が暴かれ、アインシュタインの一般相対論が新たな宇宙の法則として世界の頂点に立った歴史的瞬間でした。
\end{lecturebox}

\vspace*{3mm}
\begin{examplebox}{大人の学び直しのための実社会・重要例題：第10章}
\textbf{【実社会の具体例：1919年日食観測の熱狂と現代のダークマター探索（重力レンズ）】}\\
1919年11月、ロンドン王立協会の合同会議でアーサー・エディントン卿の日食観測結果が報告された際、会長のJ.J.トムソンは「これは人類の思想史における最も偉大な成果の一つである」と称えました。\\
(1) アインシュタインの一般相対性理論が予言した恒星光の曲がり角 $\theta = 4GM/(c^2 R)$（$= 1.75''$）が、ニュートンの重力理論に基づく計算値（$0.875''$）の正確に「2倍」になる物理的理由（時間の遅れと空間の歪みの相乗効果）を解説せよ。\\
(2) 現代天文学において、巨大銀河団による重力レンズ効果が、目に見えない「暗黒物質（ダークマター）」の質量分布を立体地図化するための最強のツールとなっている実態を述べよ。
\\[2mm]
\textbf{【懇切丁寧な解答・ステップ解説】}\\
\textbf{(1) の解：}
ニュートン力学では、空間はどこまでも平坦であり、光を仮に微小な質量を持つ粒子とみなして太陽の引力を計算すると、偏向角は $\frac{2GM}{c^2 R}$（約 $0.875$ 秒角）となります。
しかしアインシュタインの一般相対性理論では、太陽の質量によって「時間のリズム」が遅れるだけでなく、「3次元空間そのもの」も漏斗状に湾曲しています。光線はこの「曲がった時空の最短コース（測地線）」を疾走するため、空間の歪みによる寄与が正確に同量加算され、ニュートン力学のちょうど2倍である $1.75$ 秒角の偏向角が生じます。観測値はこの2倍の数値を鮮やかに実証しました。
\\[2mm]
\textbf{(2) の解：}
一般相対論の重力レンズ効果の素晴らしい特長は、「光を出さない物質であっても、質量（重力）さえあれば背後の銀河の光を曲げる」という点にあります。天文学者は、遠方銀河の歪み具合（アーク状の像や二重像）を逆算解析することで、光学望遠鏡には一切映らない謎のダークマターが宇宙のどこにどれだけ密集しているかを克明に可視化することに成功しています。
\end{examplebox}


\clearpage
\chapter{重力場の方程式とブラックホール}
\begin{center}
{\Large\bfseries\gtfamily アインシュタイン方程式の構造と極限天体の物理}
\end{center}
\vspace*{2mm}

\section*{本章の視覚的見取り図}
\begin{center}
\begin{tikzpicture}[scale=0.9, >=Stealth]
    % Flat vs Curved sheet
    \begin{scope}[yshift=2.5cm]
        \draw[thick, fill=blue!5] (0,0) -- (4,0) -- (5.5,1.5) -- (1.5,1.5) -- cycle;
        \draw[thick, red] (1.2,0.5) -- (3.8,0.5) -- (2.5,1.2) -- cycle;
        \node at (2.8, 1.8) {\textbf{平坦な空間（ユークリッド幾何学：内角の和 $= 180^\circ$）}};
    \end{scope}
    \begin{scope}[yshift=0cm]
        % Warped funnel
        \draw[thick, fill=orange!10] (0,0) .. controls (2,-1.5) and (3,-1.5) .. (5,0)
              -- (6,1.2) .. controls (3,0.2) and (2,0.2) .. (1,1.2) -- cycle;
        \shade[ball color=red!80!black] (2.7,-0.7) circle (0.4) node[below=8pt, font=\scriptsize] {大質量天体};
        \node at (3, -1.8) {\textbf{曲がった時空（リーマン幾何学：重力とは時空の歪みそのもの）}};
    \end{scope}
\end{tikzpicture}

\small\textbf{図版解説：}一般相対性理論の本質：物質とエネルギーが存在することで周囲の4次元時空連続体が歪み、重力として観測される。
\end{center}
\vspace*{2mm}

\section*{アインシュタインの原典思索}
\begin{origbox}[原著第11章：アルベルト・アインシュタインの肉声]
物質の分布と運動を表すエネルギー・運動量テンソル $T_{\mu\nu}$ は、時空の幾何学的曲率を表すアインシュタイン・テンソル $G_{\mu\nu}$ を決定する：
\[
R_{\mu\nu} - \frac{1}{2}g_{\mu\nu}R = \frac{8\pi G}{c^4} T_{\mu\nu}
\]
時空は物質にどのように動くかを指示し、物質は時空にどのように曲がるかを指示する。この方程式の最も極端な球対称解（シュワルツシルト解）は、光さえも脱出できない事象の地平面を持つ極限天体――ブラックホールの存在を必然的に導く。
\end{origbox}

\section*{大人のための徹底解説講義}
\begin{lecturebox}[徹底講義：高校数学で追体験する物理学革命]
一般相対性理論の心臓部にあるのが、1915年に完成した「アインシュタイン重力場方程式」です。
物理学者ジョン・ホイーラーはこの方程式の精神を「物質は時空に曲がり方を教え、時空は物質に動き方を教える」という名言で見事に要約しました。

アインシュタインが方程式を発表したわずか数ヶ月後、ドイツの天文学者カール・シュワルツシルトは第一次世界大戦のロシア戦線において、方程式の厳密解を計算しました。
その解が告げていたのは、天体の質量がある限界の半径（シュワルツシルト半径 $R_s = \frac{2GM}{c^2}$）よりも内側に凝縮すると、空間の歪みが無限大に達し、光の速さをもってしても二度と外へ戻れない「事象の地平面」が生まれるという衝撃的な予言でした。
当時は数学的な奇妙な解として片付けられそうになりましたが、現代では電波望遠鏡によって銀河中心のブラックホールの影が直接撮影され、アインシュタインの理論の正しさが極限の現場でも証明されています。
\end{lecturebox}

\vspace*{3mm}
\begin{examplebox}{大人の学び直しのための実社会・重要例題：第11章}
\textbf{【実社会の具体例：高層ビル展望台と地上での老化速度の違い（重力赤方偏移）】}\\
2020年、東京大学の香取秀俊教授らの研究チームは、超高精度な「光格子時計」を東京スカイツリーの地上階と展望台（高度 $450\,\text{m}$）に設置し、一般相対性理論の検証実験を行いました。\\
(1) アインシュタインの方程式が導く「強い重力場の中（地表）では時間がゆっくり進み、重力の弱い高所では時間は速く進む」という現象のメカニズムを説明せよ。\\
(2) スカイツリーの実験において、展望台の時計が地上の時計に対して1日あたり約 $4$ ナノ秒速く進んでいたという実測値が持つ、現代社会のインフラ（標高差ミリ単位の測量技術）への応用価値を述べよ。
\\[2mm]
\textbf{【懇切丁寧な解答・ステップ解説】}\\
\textbf{(1) の解：}
光が重力場を這い上がるとき、重力ポテンシャルに抗して仕事をするためエネルギーを失います。光のエネルギーは振動数に比例（$E=h\nu$）するため、エネルギーを失った光は振動数が下がり、波長が伸びます（重力赤方偏移）。時計とは「周期的な振動を数える装置」ですから、光の振動数が下がるということは、重力源に近い深い谷底にいる観測者の時計は、高い山の上にいる観測者の時計よりも「ゆっくりと時を刻む」ことを意味します。
\\[2mm]
\textbf{(2) の解：}
スカイツリーでの実験精度は $10^{-18}$（300億年に1秒の狂い）に達し、高度差わずか $450\,\text{m}$ であっても、展望台の時間が毎日約 $4.3\,\text{ナノ秒}$ 早く進むことが完全に実証されました。これは単なる物理実験を超え、「時計の進み方の差を測ることで、水準測量を行わずに地殻変動や標高差をセンチ・ミリ単位でリアルタイム監視できる」という、地震予知や防災インフラの革新的新技術として実用化が進んでいます。
\end{examplebox}


\clearpage
\part{現代宇宙論とアインシュタインの遺産}

\chapter{宇宙論と膨張する宇宙}
\begin{center}
{\Large\bfseries\gtfamily アインシュタインの生涯最大の過ちと現代ビッグバン宇宙論}
\end{center}
\vspace*{2mm}

\section*{本章の視覚的見取り図}
\begin{center}
\begin{tikzpicture}[scale=0.95, >=Stealth]
    % Expanding balloon / universe
    \draw[thick, dashed, fill=blue!5] (0,0) circle (1.2);
    \fill[red] (0.5,0.5) circle (2pt) node[right, font=\scriptsize] {銀河A};
    \fill[red] (-0.4,-0.3) circle (2pt) node[left, font=\scriptsize] {銀河B};
    \draw[thick, fill=purple!5] (4.5,0) circle (2.2);
    \fill[red] (5.4,0.9) circle (2.5pt) node[right, font=\scriptsize] {銀河A};
    \fill[red] (3.8,-0.6) circle (2.5pt) node[left, font=\scriptsize] {銀河B};
    \draw[->, ultra thick, blue] (1.5,0) -- (2.0,0) node[midway, above, font=\scriptsize] {時間発展};
    % Annotation
    \node[align=left, fill=white, draw=gray!50, rounded corners, font=\footnotesize] at (2.5,-2.8) {宇宙膨張（ハッブルの法則）：\\銀河が空間を動くのではなく、\\「空間そのものが風船のように膨張している」};
\end{tikzpicture}

\small\textbf{図版解説：}宇宙論的帰結。アインシュタイン方程式が導く動的宇宙と空間自体の膨張モデル。
\end{center}
\vspace*{2mm}

\section*{アインシュタインの原典思索}
\begin{origbox}[原著第12章：アルベルト・アインシュタインの肉声]
もし宇宙全体に物質が一様に分布しているとすれば、一般相対論の方程式は静的な平衡状態を保つことができない。重力による自己崩壊を防ぐため、私は一時的に宇宙定数（宇宙項）$\Lambda$ を方程式に導入した。
しかしハッブルによる遠方銀河の後退の発見により、宇宙そのものが膨張しているという動的宇宙論（フリードマン・ルメートル模型）が確立された。宇宙項の導入は私の生涯最大の過誤であった。
\end{origbox}

\section*{大人のための徹底解説講義}
\begin{lecturebox}[徹底講義：高校数学で追体験する物理学革命]
アインシュタインの思索の旅は、ついに「宇宙全体」の運命へと至ります。
彼が一般相対論の方程式を宇宙全体に適用したとき、方程式は恐るべき答えを突きつけてきました。「宇宙は静止していられない。膨張しているか、収縮しているかのどちらかだ」。

しかし、当時の常識は「宇宙は永遠不変で静止している」というものでした。アインシュタインは自らの方程式の予言を信じきれず、宇宙が潰れないように突っ張り棒となる反発力「宇宙項 $\Lambda$」を人工的に付け加えました。
ところが1929年、天文学者エドウィン・ハッブルが望遠鏡の観測から「遠くの銀河ほど猛烈な速度で遠ざかっている（宇宙が膨張している）」ことを発見しました。
アインシュタインは自らの修正を恥じ、「宇宙項は生涯最大の過ちだった」と撤回しました。

しかしドラマはそこで終わりませんでした。1998年、現代の天文学者は宇宙が単に膨張しているだけでなく、謎のエネルギー（ダークエネルギー）によって「加速膨張」していることを突き止めたのです。そしてその加速を説明するために復活したのが、皮肉にもアインシュタインが捨てた「宇宙項 $\Lambda$」でした。
天才の過ちさえもが、100年後の最先端天文学の鍵を握っていたのです。
\end{lecturebox}

\vspace*{3mm}
\begin{examplebox}{大人の学び直しのための実社会・重要例題：第12章}
\textbf{【実社会の具体例：宇宙背景放射（CMB）とビッグバン宇宙論、宇宙の加速膨張】}\\
テレビの地上波アナログ放送の時代、チャンネルをどの放送局にも合わない砂嵐画面にすると、「サー」という雑音（スノーノイズ）が聞こえました。その雑音の約 $1\%$ は、宇宙誕生後38万年の「宇宙の晴れ上がり」から138億年旅してきた最古の電波「宇宙マイクロ波背景放射（CMB）」でした。\\
(1) アインシュタインは当初、宇宙が永遠不変であると信じて方程式に「宇宙項 $\Lambda$」を付け加えましたが、後にハッブルの銀河後退観測を知って「生涯最大の過ち（biggest blunder）」と悔やみました。なぜ方程式は本来、動的な宇宙（膨張または収縮）を予言していたのかを説明せよ。\\
(2) 1998年の遠方超新星観測により、現代の宇宙は重力による減速ではなく「ダークエネルギーによって加速膨張している」ことが判明し、捨てられた宇宙項 $\Lambda$ が奇跡の復活を遂げた劇的なドラマを考察せよ。
\\[2mm]
\textbf{【懇切丁寧な解答・ステップ解説】}\\
\textbf{(1) の解：}
一般相対性理論の重力方程式は「物質が存在すれば空間が引き合って収縮する」ことを示しています。すべての星々が引力を持つ以上、宇宙全体が静止して安定することは数学的に不可能であり、膨張するか収縮するかの動的状態しかあり得ません。アインシュタインは当時の「静止した宇宙観」に囚われていたため、収縮を防ぐ反発力として方程式に無理やり定数（宇宙項 $\Lambda$）を挿入しました。しかし方程式自体の数学的構造は、最初から「宇宙は膨張している」というビッグバンモデルの真実を高らかに歌い上げていたのです。
\\[2mm]
\textbf{(2) の解：}
1998年、遠方のIa型超新星の精密観測により、宇宙の膨張速度が落ちるどころか、約50億年前からますます加速しているという衝撃の事実が発見されました。空間そのものを押し広げる正体不明のエネルギー「ダークエネルギー」の存在です。そして、この斥力を数学的に最もシンプルに表現するものが、アインシュタインがかつて導入し葬り去った「宇宙項 $\Lambda$」そのものでした。アインシュタインが「過ち」と呼んだものさえも、1世紀後の先端宇宙論の主役に返り咲いたのです。
\end{examplebox}


\clearpage

\appendix

\chapter{相対性理論の重要用語辞典（アルベルト・アインシュタインの語彙）}

大人の学び直しを助けるため、アインシュタインの原典および現代物理学において極めて重要な鍵概念をコンパクトに整理・解説します。

\begin{description}[font=\normalfont\bfseries\gtfamily]
    \item[慣性系（Inertial System）] 外力が働いていない物体が静止し続けるか等速直線運動を続ける座標系（慣性の法則が成立する系）。互いに等速で移動する慣性系どうしはすべて物理的に同等である。
    \item[光速度不変の原理（Principle of Constancy of Light Speed）] 真空中の光の速さ $c \approx 3 \times 10^8\,\text{m/s}$ は、光源の運動状態や観測者の運動速度に一切依存せず、すべての慣性系において常に同一であるという基本原理。
    \item[特殊相対性原理（Special Principle of Relativity）] すべての慣性座標系において、物理法則は全く同一の形式で表され、絶対静止座標系という特別な系は存在しない。
    \item[同時性の相対性（Relativity of Simultaneity）] ある観測者にとって「同時に起こった」2つの事象は、その観測者に対して動いている別の観測者にとっては「同時ではない」という事実。時間の一様性・絶対性の否定。
    \item[ローレンツ変換（Lorentz Transformation）] 光速度不変の原理と特殊相対性原理を満たす、2つの等速運動座標系間の時空変換公式。$\gamma = 1/\sqrt{1 - v^2/c^2}$ を尺度係数とする。
    \item[時間の遅れ（Time Dilation）] 観測者に対して高速運動している時計は、静止している時計に比べて進み方が遅くなる（$\Delta t = \gamma \Delta t_0$）現象。
    \item[ローレンツ収縮（Lorentz Contraction）] 観測者に対して運動している剛体は、運動方向に沿って長さが $\sqrt{1 - v^2/c^2}$ 倍に縮んで観測される現象。
    \item[質量とエネルギーの等価性（$E=mc^2$）] 静止質量 $m$ を持つ物体は、その質量自体の中に $E = mc^2$ という莫大な固有エネルギーを宿しているという帰結。原子力発電や恒星の核融合の根源。
    \item[ミンコフスキー時空（Minkowski Spacetime）] 3次元空間と時間1次元を一体として扱い、$s^2 = c^2 t^2 - (x^2 + y^2 + z^2)$ という不変量を幾何学的に定義した4次元時空連続体。
    \item[等価原理（Principle of Equivalence）] 局所的に一様な重力場の中に静止している座標系と、重力のない空間で等加速度運動をしている加速系は、物理的に全く区別がつかないという原理。
    \item[一般相対性原理（General Principle of Relativity）] 慣性系だけでなく、加速度運動している座標系や重力場の中にある座標系も含め、任意の運動状態にあるすべての座標系において物理法則が共変的（同一形式）に成り立つ。
    \item[測地線（Geodesic）] 質量やエネルギーによって歪められた曲がった時空において、物体や光線が辿る「最短（最極値）の軌道」。古典力学の「重力による引力」は、時空の測地線運動として幾何学的に書き直される。
    \item[宇宙項（Cosmological Constant, $\Lambda$）] アインシュタインが静的で閉じた宇宙解を得るために一般相対論の方程式に導入した斥力項。現代宇宙論では「ダークエネルギー」として再評価されている。
\end{description}

\chapter{総合演習問題と詳細ステップ解説（高校数学で解く相対論）}

本章では、高校数学（三平方の定理・1次方程式・平方根の展開）を用いて自らの手で計算し、相対性理論の真髄を血肉とするための総合演習問題を全10問収録しました。

\begin{exercisebox}{第1問：光時計による時間の遅れの導出}
高さ $L$ の列車内に設置された光時計を考える。床から天井へ放たれた光が反射して戻る往復時間を、車内の人 $S'$ は $\Delta t'$、堤防の静止観察者 $S$ は $\Delta t$ と測る。列車が水平方向に速度 $v$ で走るとき、三平方の定理を用いて
\[
\Delta t = \frac{\Delta t'}{\sqrt{1 - v^2/c^2}}
\]
を導出せよ。
\\[2mm]
\textbf{【懇切丁寧な解答】}\\
車内の人から見れば、光は垂直に往復するだけなので、往復距離は $2L$ である。光速を $c$ とすれば、
\[
\Delta t' = \frac{2L}{c} \quad \implies \quad L = \frac{c \Delta t'}{2}
\]
一方、堤防の観察者から見ると、光が天井に到達して戻るまでに列車は水平方向に $v \Delta t$ だけ進む。光が片道で進む水平距離は $\frac{v \Delta t}{2}$、垂直距離は $L$ である。三平方の定理より、光が進んだ斜辺の距離 $d$ は
\[
d = \sqrt{\left(\frac{v \Delta t}{2}\right)^2 + L^2}
\]
堤防の人にとっても光速は同一の $c$ であるから、往復距離 $2d$ は $c \Delta t$ に等しい：
\[
c \Delta t = 2 \sqrt{\left(\frac{v \Delta t}{2}\right)^2 + L^2} = \sqrt{(v \Delta t)^2 + 4L^2}
\]
両辺を2乗すると、
\[
c^2 (\Delta t)^2 = v^2 (\Delta t)^2 + 4L^2
\]
ここで $L = \frac{c \Delta t'}{2}$ より $4L^2 = c^2 (\Delta t')^2$ を代入すると、
\[
(c^2 - v^2) (\Delta t)^2 = c^2 (\Delta t')^2
\]
両辺を $c^2$ で割り、$(\Delta t)^2$ について解くと、
\[
\left(1 - \frac{v^2}{c^2}\right) (\Delta t)^2 = (\Delta t')^2 \quad \implies \quad \Delta t = \frac{\Delta t'}{\sqrt{1 - v^2/c^2}}
\]
これにより、高校数学の三平方の定理だけで時間の遅れの方程式が厳密に証明された。\qed
\end{exercisebox}

\vspace*{3mm}

\begin{exercisebox}{第2問：新幹線の時間の遅れの定量的計算}
東海道新幹線が時速 $288\,\text{km}$（$= 80\,\text{m/s}$）で $1$ 時間（$= 3600\,\text{秒}$）走行したとする。光速 $c = 3 \times 10^8\,\text{m/s}$ とする。\\
(1) 速度比 $\beta = v/c$ を求めよ。\\
(2) 近似式 $\sqrt{1 - x^2} \approx 1 - \frac{1}{2}x^2$ を用いて、新幹線の時計が地上の時計に対して遅れる時間をナノ秒単位で計算せよ。
\\[2mm]
\textbf{【懇切丁寧な解答】}\\
\textbf{(1) の解：}
\[
\beta = \frac{v}{c} = \frac{80}{3 \times 10^8} = \frac{8}{3} \times 10^{-7} \approx 2.67 \times 10^{-7}
\]
\textbf{(2) の解：}
地上の時間を $T = 3600\,\text{s}$、車内の時間を $T'$ とすると、$T' = T \sqrt{1 - \beta^2} \approx T \left(1 - \frac{1}{2}\beta^2\right)$ である。遅れの差 $\Delta T = T - T'$ は
\[
\Delta T \approx \frac{1}{2} T \beta^2 = \frac{1}{2} \times 3600 \times \left(\frac{8}{3} \times 10^{-7}\right)^2
= 1800 \times \frac{64}{9} \times 10^{-14} = 200 \times 64 \times 10^{-14} = 12800 \times 10^{-14} = 1.28 \times 10^{-10}\,\text{秒}
\]
これは $0.128\,\text{ナノ秒}$（$1$ ナノ秒 $= 10^{-9}$ 秒）に相当する。人間の感覚では到底認知できない微小量であるが、現代の原子時計を用いれば $10^{-18}$ 秒まで計測できるため、日常的な速度であっても相対論的遅れは完全に実証されている。\qed
\end{exercisebox}

\vspace*{3mm}

\begin{exercisebox}{第3問：素粒子ミューオンの地表到達と寿命延伸}
宇宙線が大気上層（高度 $10\,\text{km} = 10^4\,\text{m}$）で衝突して生成される素粒子ミューオンの固有寿命はわずか $\tau_0 = 2.2 \times 10^{-6}\,\text{秒}$（$2.2\,\mu\text{s}$）である。ミューオンが光速の $99.5\%$（$v = 0.995c$）で地表に向かって飛行しているとする。\\
(1) 非相対論的（ニュートン力学）な計算では、ミューオンが崩壊するまでに進める距離は何メートルか。地表に到達できるか答えよ。\\
(2) 特殊相対性理論の時間の遅れを考慮すると、地上の観測者から見たミューオンの寿命 $\tau$ は何秒に伸びるか。また、その間に進める距離を求め、なぜ地表で大量のミューオンが観測されるのかを説明せよ（ただし $\sqrt{1 - 0.995^2} \approx 0.10$ とする）。
\\[2mm]
\textbf{【懇切丁寧な解答】}\\
\textbf{(1) の解：}
古典力学での飛行距離 $L_{\text{classic}}$ は
\[
L_{\text{classic}} = v \tau_0 = (0.995 \times 3 \times 10^8\,\text{m/s}) \times (2.2 \times 10^{-6}\,\text{s}) \approx 2.985 \times 10^8 \times 2.2 \times 10^{-6} \approx 657\,\text{m}
\]
大気上層の高度は $10000\,\text{m}$ であるため、$657\,\text{m}$ 飛んだ時点で大半が崩壊してしまい、地表には一粒たりとも到達できないはずである。\\
\textbf{(2) の解：}
ローレンツ因子 $\gamma$ は
\[
\gamma = \frac{1}{\sqrt{1 - (0.995)^2}} \approx \frac{1}{0.10} = 10
\]
地上の時計から見ると、ミューオンの内部時計は $10$ 倍に引き伸ばされるため、寿命は
\[
\tau = \gamma \tau_0 = 10 \times 2.2 \times 10^{-6} = 2.2 \times 10^{-5}\,\text{秒}
\]
この間に光速に近い速度で進む距離 $L_{\text{rel}}$ は
\[
L_{\text{rel}} = v \tau \approx (3 \times 10^8) \times (2.2 \times 10^{-5}) = 6600\,\text{m} \sim 10000\,\text{m}
\]
となり、高度 $10\,\text{km}$ から放たれたミューオンの多くが地表まで余裕をもって到達できる。なお、ミューオン自身の視座（静止系）から見れば、時間の進みは $2.2\,\mu\text{s}$ のままであるが、大気の厚み $10\,\text{km}$ がローレンツ収縮によって $\frac{1}{10}$ の $1\,\text{km}$ に縮んでいるため、やはり地表に到達できると解釈され、両者の観測結果は見事に一致する。\qed
\end{exercisebox}

\vspace*{3mm}

\begin{exercisebox}{第4問：GPS衛星の相対論補正（特殊相対論と一般相対論の綱引き）}
高度 $20000\,\text{km}$ を周回するGPS衛星の軌道速度は約 $v = 3.9\,\text{km/s}$ である。\\
(1) 特殊相対性理論の速度による効果により、地上に対してGPS衛星の原子時計は1日あたり何マイクロ秒遅れるか（計算値：約 $-7.2\,\mu\text{s}$）。\\
(2) 一般相対性理論の重力ポテンシャルによる効果（高度が高く重力が弱いため時計が速く進む）により、地上に対してGPS衛星の時計は1日あたり何マイクロ秒進むか（計算値：約 $+45.8\,\mu\text{s}$）。\\
(3) 両者の総合効果として、GPS衛星の時計は地上に対して1日あたり何マイクロ秒進むか。もしこの補正を行わなければ、位置測定に1日で何キロメートルの誤差が生じるか（光速 $c = 3 \times 10^5\,\text{km/s}$ とする）。
\\[2mm]
\textbf{【懇切丁寧な解答】}\\
\textbf{(1) の解：}
1日は $86400\,\text{秒}$ である。特殊相対論の遅れは $\frac{1}{2}\frac{v^2}{c^2} \times 86400 \approx 7.2 \times 10^{-6}\,\text{秒} = 7.2\,\mu\text{s}$ の遅れ（$-7.2\,\mu\text{s}/\text{日}$）となる。\\
\textbf{(2) の解：}
高度 $20000\,\text{km}$ では地球の重力が地表よりはるかに弱いため、一般相対論により時間は速く進む。その進みは理論計算で $+45.8\,\mu\text{s}/\text{日}$ となる。\\
\textbf{(3) の解：}
正味の合成ズレは
\[
-7.2\,\mu\text{s} + 45.8\,\mu\text{s} = +38.6\,\mu\text{s}/\text{日}
\]
衛星の原子時計は地上より毎日 $38.6$ マイクロ秒早く進む。光がこの時間に進む距離は
\[
\Delta x = c \times \Delta t = (3 \times 10^5\,\text{km/s}) \times (38.6 \times 10^{-6}\,\text{s}) \approx 11.58\,\text{km}
\]
もしアインシュタインの相対性理論を考慮して送信周波数をあらかじめ微調整しておかなければ、カーナビやスマートフォンの位置情報は毎日約 $11.6\,\text{km}$ ずつズレてしまい、実用システムとして完全に崩壊する。現代社会のインフラは、相対性理論の実証によって支えられている。\qed
\end{exercisebox}

\vspace*{3mm}

\begin{exercisebox}{第5問：$E=mc^2$ と太陽の寿命・質量欠損}
太陽の放射エネルギーは毎秒約 $L_{\odot} = 3.8 \times 10^{26}\,\text{W}$（$= 3.8 \times 10^{26}\,\text{J/s}$）である。\\
(1) アインシュタインの関係式 $E = mc^2$ を用いて、太陽が光と熱として宇宙に放出することで失っている質量（質量欠損）は毎秒何トンか（光速 $c = 3.0 \times 10^8\,\text{m/s}$）。\\
(2) 太陽の現在の総質量は約 $2.0 \times 10^{30}\,\text{kg}$ である。100億年間（約 $3.15 \times 10^{17}\,\text{秒}$）このペースで燃焼し続けたとしても、失われる質量は太陽全体の何パーセント程度にすぎないかを計算し、なぜ太陽が莫大な寿命を保てるのかを述べよ。
\\[2mm]
\textbf{【懇切丁寧な解答】}\\
\textbf{(1) の解：}
1秒間に失われる質量 $\Delta m$ は
\[
\Delta m = \frac{E}{c^2} = \frac{3.8 \times 10^{26}}{(3.0 \times 10^8)^2} = \frac{3.8 \times 10^{26}}{9.0 \times 10^{16}} \approx 4.22 \times 10^9\,\text{kg} = 422万\,\text{トン}
\]
太陽は毎秒422万トンもの自身の体重をエネルギーに変換して輝いている。\\
\textbf{(2) の解：}
100億年間に失われる総質量 $M_{\text{loss}}$ は
\[
M_{\text{loss}} = 4.22 \times 10^9\,\text{kg/s} \times 3.15 \times 10^{17}\,\text{s} \approx 1.33 \times 10^{27}\,\text{kg}
\]
全質量 $M_{\odot} = 2.0 \times 10^{30}\,\text{kg}$ に対する割合は
\[
\frac{M_{\text{loss}}}{M_{\odot}} = \frac{1.33 \times 10^{27}}{2.0 \times 10^{30}} \approx 0.00067 = 0.067\%
\]
わずか $0.1\%$ 未満である。$c^2$ という係数が $9 \times 10^{16}$ という天文学的数値であるため、ごく微小な質量欠損から想像を絶する熱エネルギーが放出される。これこそが、石炭などの化学燃焼では数千年で燃え尽きてしまう太陽が、数十億年にわたって地球を温め続けられる真の理由である。\qed
\end{exercisebox}

\vspace*{3mm}

\begin{exercisebox}{第6問：相対論的速度の合成則}
時速 $100\,\text{km}$ の列車の中から前方へ時速 $50\,\text{km}$ でボールを投げると、地上から見たボールの速度はニュートン力学では単に $100 + 50 = 150\,\text{km/h}$ と足し算される。しかし相対論では速度合成則
\[
W = \frac{u + v}{1 + \frac{uv}{c^2}}
\]
に従う。\\
(1) 宇宙船が地球に対して速度 $v = 0.8c$ で飛行中、前方へ向けて速度 $u = 0.8c$ で探査機を発射した。地球から見た探査機の速度 $W$ は光速の何倍か。光速 $c$ を超えないことを示せ。\\
(2) 宇宙船から前方へ光（$u = c$）を発射した場合、地球から見た光の速度 $W$ を合成則から計算し、光速度不変の原理と矛盾しないことを確かめよ。
\\[2mm]
\textbf{【懇切丁寧な解答】}\\
\textbf{(1) の解：}
\[
W = \frac{0.8c + 0.8c}{1 + \frac{(0.8c)(0.8c)}{c^2}} = \frac{1.6c}{1 + 0.64} = \frac{1.6c}{1.64} \approx 0.9756c
\]
古典力学の足し算なら $1.6c$ となり光速を突破するが、分母の補正項 $1 + uv/c^2$ の働きにより、結果は光速の約 $97.6\%$ に留まり、決して光速を超えない。\\
\textbf{(2) の解：}
$u = c$ を代入すると、
\[
W = \frac{c + v}{1 + \frac{cv}{c^2}} = \frac{c + v}{1 + \frac{v}{c}} = \frac{c(c + v)}{c + v} = c
\]
どれほど高速（$v$）で運動している乗り物から光を放っても、外の観測者から見た光の速さは寸分違わず $c$ である。速度合成則自体の中に光速度不変の公理が完璧に織り込まれている。\qed
\end{exercisebox}

\vspace*{3mm}

\begin{exercisebox}{第7問：等価原理とエレベーターの光の曲がり}
無重力空間を上向きに一定加速度 $g$ でロケット（またはエレベーター）が上昇している。幅を $W$ とする。左側の壁の小穴から水平に入射した光線が、右側の壁に到達する現象を考える。\\
(1) 光が右側の壁に到達するまでの時間 $t$ を光速 $c$ と $W$ を用いて表せ。\\
(2) その時間 $t$ の間にエレベーターが上向きに加速することによって上昇する距離 $h = \frac{1}{2}gt^2$ を求めよ。\\
(3) エレベーター内の観測者から見ると、光線が下向きに放物線を描いて $h$ だけ落ちて見えることを説明し、アインシュタインの等価原理により「重力場の中でも光線は必ず下方に湾曲する」と予言できる論理を述べよ。
\\[2mm]
\textbf{【懇切丁寧な解答】}\\
\textbf{(1) の解：}
光の水平速度は $c$ であるから、幅 $W$ を横切る時間は $t = \frac{W}{c}$ である。\\
\textbf{(2) の解：}
等加速度運動の公式より、上昇距離は
\[
h = \frac{1}{2} g t^2 = \frac{g W^2}{2 c^2}
\]
\textbf{(3) の解：}
外の慣性系から見れば光は真横に直進しているが、床が上に向かって加速して迎え撃つため、車内の乗客から見れば光線は床に向かって $h$ だけ下方に曲がったように観測される。
等価原理によれば「加速度系と重力場は物理的に同一」であるから、地上に静止したエレベーター（重力加速度 $g$）の中でも、全く同一の湾曲 $h$ が光線に生じなければならない。こうしてアインシュタインは、「質量のない純粋な光線も重力によって曲げられる」という驚天動地の結論を、高校物理の運動方程式だけで導き出したのである。\qed
\end{exercisebox}

\vspace*{3mm}

\begin{exercisebox}{第8問：太陽による恒星光の重力偏向角（エディントン観測の検証）}
1919年5月29日の皆既日食において、アーサー・エディントン卿は西アフリカのプリンシペ島で太陽の縁すれすれを通る恒星の光を観測した。一般相対性理論による光の偏向角の公式は
\[
\theta = \frac{4GM}{c^2 R} \quad \text{（ラジアン）}
\]
で与えられる。\\
太陽質量 $M_{\odot} \approx 1.99 \times 10^{30}\,\text{kg}$、太陽半径 $R_{\odot} \approx 6.96 \times 10^8\,\text{m}$、万有引力定数 $G \approx 6.67 \times 10^{-11}\,\text{m}^3/(\text{kg}\cdot\text{s}^2)$、光速 $c \approx 3.00 \times 10^8\,\text{m/s}$ とする。\\
(1) 偏向角 $\theta$ をラジアン単位で求めよ。\\
(2) $1\,\text{ラジアン} = \frac{180 \times 3600}{\pi} \approx 206265\,\text{秒角}$（arcsec）を用いて、偏向角を秒角に換算せよ。古典ニュートン力学（光を質量を持つ粒子と仮定した場合の半分）の予言値 $0.875''$ と比較せよ。
\\[2mm]
\textbf{【懇切丁寧な解答】}\\
\textbf{(1) の解：}
\[
\theta = \frac{4 \times (6.674 \times 10^{-11}) \times (1.989 \times 10^{30})}{(3.00 \times 10^8)^2 \times (6.963 \times 10^8)}
= \frac{5.310 \times 10^{20}}{6.267 \times 10^{25}} \approx 8.473 \times 10^{-6}\,\text{rad}
\]
\textbf{(2) の解：}
\[
\theta'' = 8.473 \times 10^{-6} \times 206265'' \approx 1.75''
\]
ニュートン力学の粒子論による計算では $\frac{2GM}{c^2 R} = 0.875''$ であったのに対し、アインシュタインの一般相対性理論では「時間の遅れ」と「空間の湾曲」の双方が寄与するため、正確に2倍の $1.75''$ となる。エディントンの観測隊が日食写真の乾板を測定した結果は $1.98'' \pm 0.12''$ および $1.61'' \pm 0.30''$ であり、アインシュタインの予言が鮮やかに立証され、世界中の新聞が「科学の歴史が書き換わった」と一面で報じた。\qed
\end{exercisebox}

\vspace*{3mm}

\begin{exercisebox}{第9問：水星の近日点移動とニュートン力学の破綻}
太陽に最も近い惑星である水星の楕円軌道は、100年間に約 $5600''$（秒角）回転（移動）する。そのうち木星や金星など他の惑星からの摂動によるニュートン力学的なズレは $5160''$、地球の歳差運動による見かけの回転が $43''$ 分解され、どうしても説明のつかない「余剰の $43''$/世紀」が19世紀以来の天文学の最大の謎であった。\\
アインシュタインが導いた一般相対論による1公転あたりの近日点移動角 $\Delta \phi$ は
\[
\Delta \phi = \frac{6\pi GM}{c^2 a (1 - e^2)} \quad \text{（ラジアン/公転）}
\]
で与えられる。水星の軌道長半径 $a \approx 5.79 \times 10^{10}\,\text{m}$、離心率 $e \approx 0.2056$、公転周期 $T \approx 87.97\,\text{日}$（1世紀あたりの公転回数 $N \approx 415$ 回）として、1世紀あたりの相対論的移動角を秒角で計算せよ。
\\[2mm]
\textbf{【懇切丁寧な解答】}\\
分母の幾何学的項を計算すると、
\[
1 - e^2 = 1 - (0.2056)^2 = 1 - 0.0423 = 0.9577
\]
1公転あたりの角移動 $\Delta \phi$ は
\[
\Delta \phi = \frac{6\pi \times (6.674 \times 10^{-11}) \times (1.989 \times 10^{30})}{(3.00 \times 10^8)^2 \times (5.79 \times 10^{10}) \times 0.9577}
= \frac{2.499 \times 10^{21}}{4.991 \times 10^{27}} \approx 5.007 \times 10^{-7}\,\text{rad/公転}
\]
1世紀あたりの総移動角は、公転回数 $N = 415$ を掛けて、
\[
\Delta \Phi = 415 \times 5.007 \times 10^{-7} \approx 2.078 \times 10^{-4}\,\text{rad/世紀}
\]
これを秒角に換算すると、
\[
\Delta \Phi'' = 2.078 \times 10^{-4} \times 206265'' \approx 42.86'' \approx 43''/\text{世紀}
\]
1世紀にわずか $43''$（満月の視直径の約 $\frac{1}{40}$）という極微の狂いが、アインシュタインの方程式によって何ら任意の仮定なしにピタリと導き出された。この計算に成功した夜、アインシュタインは「心臓が激しく波打ち、全身が震えるほどの歓喜を覚えた」と友人に書き送っている。\qed
\end{exercisebox}

\vspace*{3mm}

\begin{exercisebox}{第10問：シュワルツシルト半径（ブラックホールの境界面）}
一般相対性理論において、天体の質量 $M$ がそのシュワルツシルト半径
\[
R_s = \frac{2GM}{c^2}
\]
の内側にまで凝縮すると、表面脱出速度が光速 $c$ に達し、光さえも脱出できない「事象の地平面（ブラックホール）」が形成される。\\
(1) 地球の質量 $M_{\text{earth}} \approx 5.97 \times 10^{24}\,\text{kg}$ をブラックホールにするには、何ミリメートルの半径に圧縮しなければならないか。\\
(2) 太陽の質量 $M_{\odot} \approx 1.99 \times 10^{30}\,\text{kg}$ のシュワルツシルト半径は何キロメートルか。\\
(3) 銀河系の中心にある巨大ブラックホール「いて座A*（Sagittarius A*）」の質量は太陽の約400万倍である。そのシュワルツシルト半径は何万キロメートルか。地球と太陽の距離（1天文単位 $\approx 1.5 \times 10^8\,\text{km}$）と比較せよ。
\\[2mm]
\textbf{【懇切丁寧な解答】}\\
\textbf{(1) の解：}
\[
R_s = \frac{2 \times (6.674 \times 10^{-11}) \times (5.972 \times 10^{24})}{(3.00 \times 10^8)^2} = \frac{7.971 \times 10^{14}}{9.00 \times 10^{16}} \approx 8.86 \times 10^{-3}\,\text{m} \approx 8.9\,\text{mm}
\]
地球の全物質を、わずかビー玉サイズ（直径約 $1.8\,\text{cm}$）にまで押し縮めるとブラックホールになる。\\
\textbf{(2) の解：}
\[
R_s = \frac{2 \times (6.674 \times 10^{-11}) \times (1.989 \times 10^{30})}{(3.00 \times 10^8)^2} = \frac{2.655 \times 10^{20}}{9.00 \times 10^{16}} \approx 2950\,\text{m} \approx 2.95\,\text{km} \approx 3\,\text{km}
\]
太陽ほどの巨大な恒星でも、半径約 $3\,\text{km}$（富士山の高さ程度）に収縮すればブラックホールとなる。\\
\textbf{(3) の解：}
太陽の400万倍であるから、
\[
R_s = 4.0 \times 10^6 \times 2.95\,\text{km} \approx 1.18 \times 10^7\,\text{km} = 1180万\,\text{km}
\]
太陽の半径が約 $70万\,\text{km}$ であるから、その約17倍の領域である。地球と太陽の距離 $1億5000万\,\text{km}$ の約 $8\%$（水星の軌道半径の約 $\frac{1}{5}$）にすっぽりと収まるサイズである。2022年にイベント・ホライズン・テレスコープ（EHT）が撮影したいて座A*の電波写真の影（シャドウ）は、まさにこのアインシュタイン方程式の解が1世紀の時を経て天体観測によって視覚的に捉えられた記念碑的成果である。\qed
\end{exercisebox}

\chapter*{おわりに：現代文明の礎となったアインシュタインの遺産}
\addcontentsline{toc}{chapter}{おわりに：現代文明の礎となったアインシュタインの遺産}

アルベルト・アインシュタインが本書『特殊および一般相対性理論について』を出版してから一世紀以上の歳月が流れました。かつて「世界で真に理解できる者は3人しかいない」とまで噂されたこの理論は、いまやSFの空想物語ではなく、私たちのスマートフォンの中（GPS高精度測位）、病院の先端医療（PET診断における陽電子消滅）、エネルギーインフラ（原子力発電）、そして宇宙望遠鏡による深宇宙探査の現場において、寸秒の狂いもなく稼働する「現代文明の実用技術」となっています。

しかし、技術的応用の価値以上に私たちの心を揺さぶり続けるのは、アインシュタインが見せた\textbf{「人間の理性と論理的思考の力」}です。彼は実験室に籠もって大がかりな装置を操作したのではなく、自らの頭の中で「もし光の速さで光を追いかけたらどう見えるか」「走る列車の中で雷を見たらどうなるか」「落下するエレベーターの中の人には何が起きるか」という純粋な\textbf{思考実験（Gedankenexperiment）}を徹底的に突き詰めることで、2000年来の絶対時間・絶対空間の常識を覆しました。

大人の学び直しの旅を終えた読者の皆様が、夜空を見上げるとき、あるいは手元のスマートフォンの地図アプリを開くとき、その背後に広がる「時間と空間が織りなす壮大で美しい四次元幾何学の世界」を、アインシュタインの温かな眼差しとともに実感していただけるなら、編訳者一同のこの上ない喜びです。

\end{document}
